%

\chapter{{Korpusstudie}} \label{Korpus}

In diesem Kapitel wird die Korpusanalyse präsentiert, die durchgeführt wurde, um eine Stichprobe an Intensitätspartikeln quantitativ einzuordnen.\footnote{Auszüge dieses Kapitels wurden in \citet[]{Schmidt.2022} veröffentlicht.} Wie in Abschnitt \ref{Funktionsklasse} erläutert, wird in der theoretischen Literatur die Position vertreten, dass Intensivierer eine Verschiebung von Graden auf einer Skala bewirken, wodurch dem jeweiligen Intensivausdruck eine stärkere Intensität und somit eine höhere Skalenposition zukommt. Der Unterschied zwischen deskriptiven und expressiven Intensivierern liegt darin, dass letztere mutmaßlich nicht nur skalierend sind, sondern auch die persönliche Einstellung des Sprechers zur geäußerten Sachlage ausdrücken (vgl. \citealt[315]{Lang.1983}; \citealt[57]{Ortner.2014}). Das Ziel dieser Untersuchung ist es, auf Grundlage eines Zeitschriftenkorpus die Frequenzverläufe verschiedener expressiver Intensivierer des Gegenwartsdeutschen (n = 16) über einen Zeitraum von rund 70 Jahren zu erfassen. Dadurch wird ermittelt, wie sich die Intensiviererklassen synchron und diachron quantitativ in einem schriftsprachlichen Kontext verhalten. Mithilfe der Ergebnisse können erste Rückschlüsse auf die Expressivität der Ausdrücke gezogen und Annahmen für weiterführende experimentelle Untersuchungen abgeleitet werden.

\section{Theoretischer Rahmen und Forschungshypothesen} \label{Forschungshypothesen-Korpus}

Wie einleitend geschildert, basiert die Untersuchung auf der Beobachtung, dass es im Deutschen eine Vielzahl von Intensivierern gibt. Dies wirft die Frage auf, welche Gründe es für den Variantenreichtum gibt. In diesem Kontext wird in der Literatur traditionell die These vertreten, dass diese Partikeln einem schnellen und unabdingbaren Wandel unterliegen (vgl. \citealt[2]{Hauschild.1899}; \citealt[267]{Becher.1907}; \citealt[274]{Baumgarten.1908}; \citealt[126f.]{Biedermann.1969};; \citealt[18]{Bolinger.1972}; \citealt[216f.]{Keller.1995};; \citealt[257]{Ito.2003}).\footnote{Neben dieser Erklärung gibt es in der soziolinguistischen Literatur auch soziale Ansätze zur Erklärung der Variation, indizieren Intensivierer doch auch soziale Bedeutung. So liefert der von einem Sprecher gewählte Ausdruck indirekt auch Informationen über Herkunft, Alter oder Geschlecht der Person. Auch wenn der Schwerpunkt dieses Buches nicht auf den sozialen Faktoren liegt, die zur Variation von Intensivierern beitragen, ist die soziale Dimension dennoch eng mit dem Phänomen der Intensivierungsvariation verknüpft.} Ihre Veränderlichkeit lässt sich laut \citet[189ff.]{Kirschbaum.2002b} primär auf das Streben nach Expressivität und Originalität zurückführen, durch das der Sprecher die größtmögliche kommunikative Wirkung beim Adressaten erzielen möchte: \glqq{}Wer in bestimmten Gruppen besonders expressive Ausdrücke verwendet, `kommt gut an'\grqq{} (\citealt[3]{Keller.2003}). Ist die Annahme, dass Entwicklungen von Intensivierern schnell vonstatten gehen und daher leicht nachzuvollziehen sind, faktisch gültig, so ist davon auszugehen, dass sie die ideale Basis für diachrone Korpusanalysen bieten.

Der Hauptantrieb für den Wandel des lexikalischen Inventars wird im Allgemeinen in zu hoher Frequenz gesehen, die ursächlich für den schnellen Verschleiß expressiver Ausdrücke sei und letztlich dazu führe, dass die intendierte Wirkung nachlasse: \glqq{}[W]eil jeder gerne imponieren möchte, werden solche Ausdrücke häufig verwendet; und wenn sie häufig verwendet werden, verlieren sie ihren expressiven `Pfiff'.\grqq{} (\citealt[3]{Keller.2003}). Mit anderen Worten: Je häufiger ein Intensivierer gebraucht wird, desto vorhersehbarer und weniger originell ist er. Dies wirkt sich in der Konsequenz auf seine Expressivität aus, die nach Ansicht von \citet[2]{Hauschild.1899} unweigerlich schwindet. Dadurch wird laut \citet[274]{Baumgarten.1908} neuen, expressiveren Ausdrücken der Weg in die Intensiviererrolle geebnet.\footnote{Als Beispiel für einen quantitativ rückläufigen Intensivierer führt \citet[10f.]{Hauschild.1899} \textit{wunder} an, das bereits im Althochdeutschen existierte, im Mittelhochdeutschen inflationär und seinerzeit nur noch selten gebraucht worden sei. Heute findet man \textit{wunder} bloß noch als Präfixoid in festen Konstruktionen wie bspw. \textit{wunderschön} (vgl. auch \citealt[228]{Pittner.1989}).}
Diese Gedankenfolge bildet den Ausgangspunkt dieser explorativen Korpusstudie, die den Entwicklungsgang von Intensivierern über die Zeit untersucht. Dadurch wird die Gültigkeit der folgenden häufig vertretenen, jedoch nicht empirisch abgesicherten Hypothesen überprüft:

\begin{itemize}
\item[\textbf{\textsc{H1}}] Eine hochfrequente Verwendung von Intensivierern führt unweigerlich zum Verlust ihrer Expressivität.

\item[\textbf{\textsc{H2}}] Der Expressivitätsverlust geht mit einem Bedarf an neuen Ausdrücken einher, mit denen sich Sprecher auch weiterhin als expressiv, innovativ und originell darstellen können.
\end{itemize}

\noindent Dabei lässt vor allem die Hypothese H1 Raum für Spekulation. In der Literatur wird gemeinhin als unkontrovers betrachtet, dass ein übermäßiger Gebrauch von expressiven Ausdrücken in unerwünschter Abnutzung und unvermeidlichem Expressivitätsverlust resultiert. Was die Hypothese nach meinem Verständnis allerdings ebenso zu beinhalten scheint, ist der Faktor Zeit. So ist anzunehmen, dass sich der Expressivitätsstatus eines Intensivierers nicht abrupt von expressiv zu nicht-expressiv verändert, sondern einem graduellen Wandel unterliegt. Dabei spielt vermutlich inhärent auch die Fortdauer seines Gebrauchs eine zentrale Rolle -- ein Aspekt, der meines Wissens bisher nicht berücksichtigt wurde. Dennoch hebt schon \citet[]{Tobler.1868} den Einfluss von Zeit wie folgt hervor:

\begin{singlespace*}
\begin{quotation} \noindent Im lauf der zeit versinkt dann die qualitative grundlage der bedeutung des ersten wortes immer mehr, und es bleibt blo\textlongs s gleichsam ein quantitativer niederschlag derselben, ein allgemeiner begriff von steigerung des zweiten durch das erste als seinen exponenten. \hfill (\citealt[59]{Tobler.1868}) \\ \end{quotation}
\end{singlespace*}

\noindent Abgesehen von der Verwendungshäufigkeit kommt also auch ein zweiter Parameter als Triebkraft für Expressivitätsverlust infrage, nämlich die Zeit. Um dem gerecht zu werden, wird die Hypothese H1 gemäß der soeben dargelegten Faktoren Frequenz und Zeit folgendermaßen spezifiziert: 

\begin{itemize} 
\setlength{\itemindent}{1cm}
\item[\textbf{\textsc{H1-Freq}}] Eine hochfrequente Verwendung von Intensivierern führt zu einer pragmatischen Abnutzung und demzufolge sukzessive zum Verlust von Expressivität. 

\item[\textbf{\textsc{H1-Temp}}] Eine persistente Verwendung von Intensivierern führt zu einer pragmatischen Abnutzung und demzufolge sukzessive zum Verlust von Expressivität. 
\end{itemize}

Da mittels einer frequenzbasierten Korpusuntersuchung lediglich das diachrone Vorkommen von Ausdrücken in einer Textsammlung erhoben wird, erlauben die ermittelten Daten naturgemäß keine Aussagen über den Expressivitätsgrad der Intensivierer. Nichtsdestotrotz liefern sie erste Informationen über die Rolle von Frequenz und Persistenz und können als Spiegelbild des potenziellen Expressivitätsverlusts dienen. Die longitudinale Entwicklung wird daher als erster Hinweis auf den Expressivitätsgrad eines jeden Intensivierers herangezogen. Im Hinblick auf die Hypothese H1 wird im Sinne der in Abschnitt \ref{Auslöser} beschriebenen Expressivitätsmaximen nach \citet[]{Kirschbaum.2002b} erwartet, dass die Frequenz eines neuen expressiven Intensivierers zunächst ansteigt, da er bei Sprechern aufgrund seiner bemerkenswerten Wirkung beliebt ist, er mit der Zeit und durch übermäßigen Gebrauch allerdings an Expressivität einbüßt. Dadurch wird er im weiteren Zeitverlauf entweder seltener (bspw. \textit{wunder})\footnote{In diesem Zusammenhang betont \citet[52]{Gevaudan.2007}, dass lexikalisierte Ausdrücke ohnehin nicht notwendigerweise persistent seien, sondern \glqq{}immer wieder gebraucht werden [müssen], sonst verfallen sie. \textelp{} Wenn die Verwendungshäufigkeit für ein Wort abnimmt, bis sie unter eine bestimmte Schwelle fällt, führt das zu lexikalischem Schwund.\grqq{}} oder mit rein skalierender Funktion (bspw. \textit{sehr}) verwendet. Letzteres hält \citet[12]{Hauschild.1899} zum intensivierenden \textit{blut} wie in \textit{blutjung} fest, das sich etwa im 17. Jahrhundert etabliert hätte, später jedoch seinen semantischen Gehalt verloren habe und infolgedessen \glqq{}nur noch als Verstärkungswörtchen gefühlt\grqq{} worden sei. Die Frequenzabnahme eines Ausdrucks wird mutmaßlich von einer Zunahme eines anderen, expressiveren Intensivierers begleitet. Durch den neuen Intensivierer können Sprecher den Verlust eines abgenutzten Ausdrucks kompensieren und auch zukünftig expressiv bleiben. Diese Vermutung geht auf \citet[]{Baumgarten.1908} zurück, der sich über Steigerungskomposita wie \textit{blitzblank} wie folgt äußert: \glqq{}So entstehen \textelp{} Wörter, an denen das Empfindungsleben in höchstem Maße beteiligt ist, die aber auch hier und da wieder verblassen, durch neue ersetzt oder in derselben Weise doppelt und dreifach gesteigert werden.\grqq{} (\citealt[274]{Baumgarten.1908}). Bezüglich der Hypothese H2 wird also erwartet, dass innerhalb des Untersuchungszeitraums neue expressive Intensivierer auftreten,\footnote{\textit{Untersuchungszeitraum} meint hier und im Folgenden die Zeitspanne, über die sich das Korpus erstreckt, und nicht den Zeitraum, in dem die Korpusanalyse durchgeführt wurde.} die ein Indiz für einen tatsächlich bestehenden Bedarf an neuen Ausdrücken sind und deren Gebrauchsfrequenz, im Einklang mit H1, ebenso vorerst zunimmt. Ein Zutreffen von H2 würde sich folglich in Form von neuen Ausdrücken niederschlagen, die sich erst im Laufe der Jahre erstmals in intensivierender Funktion im Korpus nachweisen lassen.

\largerpage
Zusammengenommen liegt das Ziel der Untersuchung darin, die bis dato als unzweifelhaft erachteten Hypothesen erstmals empirisch zu überprüfen, um stichhaltige Aussagen über deren Gültigkeit zu machen. Die erste stammt von \citet[2]{Hauschild.1899},\footnote{In der Literatur wird an dieser Stelle meist auf \citet[126f.]{Biedermann.1969}, \citet[99]{Suscinskij.1985} oder \citet[2]{Keller.2003} verwiesen. Dessen ungeachtet ist \citet[2]{Hauschild.1899}, soweit mir bekannt, der Erste, der die große Varianz von Intensitätsausdrücken bzw. Verstärkungskomposita auf die Frequenz als kausale Ursache zurückführt.} der anmerkt, dass Intensivierer zwangsläufig ihren expressiven Charakter verlieren. Diesen Verlust sehen er und andere Autor:innen hauptsächlich mit der Häufigkeit des Gebrauchs assoziiert: \glqq{}Frequenz ist der natürliche Feind von Expressivität.\grqq{} (\citealt[2]{Keller.2003}). Wie \citet[59]{Tobler.1868} und \citet[253]{Gabelentz.2016} beschreiben, hat neben der Frequenz eventuell auch die Dauer des Gebrauchs einen entscheidenden Einfluss auf die Expressivität. Die damit verbundene Abnutzung setze einen nicht-intendierten innovativen Wandel in Gange, der sich nach Auffassung von \citet[274]{Baumgarten.1908} durch neue Ausdrücke bemerkbar macht.\footnote{Wie schon \citet[]{Hauschild.1899} findet auch \citet[]{Baumgarten.1908} in neueren Arbeiten keine Erwähnung, sodass diesbezüglich im Allgemeinen auf Autor:innen wie z.\,B. \citet[18]{Bolinger.1972}, \citet[99]{Suscinskij.1985} oder \citet[391]{Tagliamonte.2008} Bezug genommen wird. Weil \citet[274]{Baumgarten.1908} sich dem Phänomen erheblich früher widmet, beziehe ich mich in diesem Buch folgerichtig auf ihn.} Diese ermöglichen Sprechern, die Expressivität trotz des Wegfalls altgedienter Formen beizubehalten, was zugleich die zweite der zu testenden Hypothesen ist.

\section{Methodisches Vorgehen} \label{Durchführung}

Die empirische Grundlage bildet eine korpuslinguistische Analyse von synchronen und diachronen Häufigkeitsverteilungen ausgewählter Intensivierer des Gegenwartsdeutschen. Durch dieses Vorgehen können mögliche Wandelprozesse, die sich im (in dieser Untersuchung schriftlichen) Sprachgebrauch abzeichnen, sichtbar und interpretierbar gemacht werden. Dabei besteht die Möglichkeit, dass Veränderungen in der Schriftsprache Rückschlüsse auf die gesprochene Sprache zulassen. So ist vorstellbar, dass der schriftliche Sprachgebrauch mitunter vom mündlichen Sprachgebrauch beeinflusst ist -- jedenfalls sofern es stilistisch der jeweiligen Textsorte entspricht. Wenn es die Textsorte gestattet, kann der schriftliche Sprachgebrauch gewissermaßen als Abbild des mündlichen Sprachgebrauchs betrachtet werden \citep[vgl.][]{Koch.1996,Koch.2007}.\footnote{Nichtsdestotrotz wird ausdrücklich hervorgehoben, dass es sich hierbei um schriftliche und nicht um gesprochene Daten handelt.} Zudem erlaubt die quantitative Einordnung der untersuchten Ausdrücke erste Mutmaßungen in Bezug auf die Rolle der Expressivität sowie den zu überprüfenden Expressivitätsverlust.

Korpusstudien haben im Vergleich zu Experimenten den Vorteil, dass durch beträchtliche Textsammlungen authentische, d.\,h. natürlichsprachliche Daten, zur %quantitativen und qualitativen
Analyse herangezogen werden können. Diese sind in elektronischer Form im Netz größtenteils frei zugänglich und bestenfalls bereits digital aufbereitet. Da Intensitätspartikeln dem Anschein nach zu schnellem Wandel neigen (vgl. \citealt[257]{Ito.2003}; \citealt[281]{Tagliamonte.2005}), bieten sie eine hervorragende Grundlage für Korpusanalysen, mithilfe derer synchrone und diachrone Häufigkeitsverteilungen ermittelt und Annahmen über die Rolle der Expressivität abgeleitet werden können. Die dadurch gewonnenen Daten können außerdem als Basis für weiterführende experimentelle Untersuchungen dienen.

\subsection{Belegkorpus}
Das als Stichprobe herangezogene Untersuchungskorpus wurde über das Deutsche Referenzkorpus [\textsc{DeReKo}] aus dem \textit{W\,-\,Archiv der geschriebenen Sprache} bzw. \textit{W\,-\,öffentlich} durch das Webinterface der Korpusmanagement- und Recherchesoftware \textsc{COSMAS II} (Version 2.3.5)\footnote{\citet[]{COSMAS.2020}: \textit{https://cosmas2.ids-mannheim.de/cosmas2-web/} \mbox{[letzter Aufruf: 25.02.2020]}.} extrahiert \citep[vgl.][]{Kupietz.2010}. Das \textsc{DeReKo} enthält über 350 elektronische Korpora geschriebener Sprache mit mehr als 42 Milliarden Textwörtern (Stand: 06.02.2018) und wird vom Leibniz-Institut für Deutsche Sprache in Mannheim (\citealt[]{IDS.2018}) kostenfrei bereitgestellt. Dessen Grundlage bilden alle systematisch konzipierten und öffentlich zugänglichen Textkorpora des W-Archivs (mit Neuakquisitionen), die allesamt morphosyntaktisch annotiert sind, insgesamt etwa 32 Millionen deutschsprachige Texte mit neun Milliarden Wortformen umfassen und somit vermutlich hinreichend repräsentative Belege für linguistische Analysen bereithalten. Das zurate gezogene \textit{W\,-\,Archiv der geschriebenen Sprache} setzt sich vornehmlich aus Zeitungs- bzw. Zeitschriftentexten, aber auch Protokollen aus Parlamentsdebatten sowie Beiträgen aus Diskussionsforen zusammen. Auf dieser Grundlage wurde ein eigenes virtuelles Untersuchungskorpus definiert, das alle Texte der über das \textsc{DeReKo} zur Verfügung gestellten digitalen Ausgaben der überregionalen, wöchentlich erscheinenden Zeitschriften \textsc{Der Spiegel}\footnote{Streng genommen handelt es sich beim \textsc{Spiegel} um ein Nachrichtenmagazin; er wird im Kontext dieser Untersuchung jedoch gemeinsam mit der \textsc{Zeit} als Zeitschrift bezeichnet.} im Zeitraum von Januar 1950 bis Dezember 2017 sowie \textsc{Die Zeit} im Zeitraum von Januar 1953 bis Dezember 2017 beinhaltet.

Diese beiden Zeitschriften wurden als Untersuchungsgrundlage gewählt, weil sie es ermöglichen, ein offenbar durch schnellen Wandel geprägtes Sprachphänomen in einer homogenen Publikationsform \citep[vgl.][61]{Adamzik.2016} über eine große Zeitspanne hinweg zu inspizieren. Vor allem letztere ist vonnöten, um eine umfassende linguistische Auswertung auf Basis relevanter gegenwartssprachlicher Daten durchzuführen. Daher wird angenommen, dass das zusammengestellte Zeitschriftenkorpus eine geeignete Quelle für die Untersuchung von sprachlicher Intensivierung darstellt. Die ermittelten Gebrauchsfrequenzen lassen idealerweise erste Rückschlüsse auf die Gültigkeit der oben angeführten Hypothesen H1 und H2 nach \citet[2]{Hauschild.1899} und \citet[274]{Baumgarten.1908} zu.

Dennoch gilt es zu beachten, dass es sich bei der eher konservativen Zeitungssprache, dem Modell der Einteilung von Interaktionsformen nach \citet[]{Koch.1996,Koch.2007} folgend, um konzeptionell schriftlichen Sprachgebrauch handelt, der sich durch die kommunikative Distanz der (geschriebenen) Standardsprache \citep[vgl.][217,\,226]{Eisenberg.2007} auszeichnet. Aufgrund der Einhaltung sprachnormativer Regeln kann es darin mitunter schwierig sein, ausreichend Belege für informelle umgangs- und gruppensprachliche Ausdrücke -- insbesondere neue, hochaktuelle Intensivierer wie \textit{krass} oder \textit{übelst} -- zu finden, denn \glqq{}[a]ls natürliche Kraft \textelp{} tritt Expressivität in der Nähekommunikation besonders hervor \textelp{}, dagegen wird sie in der Distanz meist unterdrückt.\grqq{} (\citealt[12]{Pustka.2014b}). % Gleiches berichtet auch \citet[19]{Berz.1953}: "Der T. st. (= Typus steinreich, Anmerkung J. S.) als eine durch emotionale Triebkräfte entstandene Fehlleistung eines Individuums hat es nicht leicht, sich zum allgemeinen Sprachgebrauch durchzusetzen. Er dringt in der Regel nicht in die Schriftsprache, weil er der begrifflich-logischen Zensur nicht standhält und weil für ihn dort die Lebensbedingungen zu karg sind."
Gleiches berichtet \citet[12]{Helbig.1994}, der den Unterschied zwischen partikelarmer Schriftsprache und partikelreicher Umgangs- oder Alltagssprache nicht bloß auf Texttypen im Allgemeinen zurückführt. Stattdessen sehe er ihn im Besonderen darin, \glqq{}ob sich der Sprecher direkt an einen präsenten Partner (im Dialog) oder nur per Distanz an einen weder präsenten noch spezifizierten Partner wendet. \textelp{} Je umgangssprachlicher ein Text ist, desto partikelreicher ist er in der Regel auch.\grqq{} In diesem Zusammenhang verweist er auf \citet[98ff.]{Weydt.1969}, dessen Ergebnissen zufolge die Zeitungssprache quantitativ gerade einmal über ein Drittel der Partikeln aus der Alltagssprache verfügt. Daneben kommt auch \citet[451]{Darcy.2015} zu dem Schluss, dass Intensivierung eher im mündlichen als im schriftlichen Diskurs auftritt. Außerdem hat auch die Forschung gezeigt, dass der Gebrauch von Intensivierern in narrativen Texten häufiger ist als in nicht-narrativen Texten und dass auch das Register die Variation von Intensivausdrücken beeinflusst (vgl. u.\,a. \citealt[8]{Brown.2012}). Ebenso ist zu erwarten, dass mögliche Veränderungen im Vergleich zur gesprochenen Kommunikation, die gemeinhin als Motor für expressive Innovationen gilt, im Korpus zeitverzögert auftreten. So müssen sich neue Ausdrücke erst einmal den Weg in solche speziellen Textformen bahnen, was nicht ad hoc geschieht, sondern verständlicherweise einige Zeit in Anspruch nimmt. Dies birgt zudem die Gefahr, dass die Intensivierer in der Zwischenzeit im mündlichen Sprachgebrauch wieder obsolet sind und es ihnen erst gar nicht gelingt, in die Schriftsprache überführt zu werden. Weiterhin ist davon auszugehen, dass vor allem die neueren expressiven Intensivierer erheblich seltener im Korpus vertreten sind als ältere, bereits etablierte Ausdrücke, gelten expressive Intensivierer doch ohnehin als Eigenart informeller, mündlicher Substandards (vgl. \citealt[143]{Lorenz.2002}; \citealt[362]{Tagliamonte.2008}). So legt bspw. \citet[40]{Paradis.1997} im Rahmen ihrer korpusgestützten Studie zu englischen Intensivierern dar, dass die von ihr überprüften Ausdrücke fast ausschließlich im mündlichen Korpus auftreten, während lediglich einige wenige auf das schriftliche Korpus beschränkt zu sein scheinen. Gleichsam stellt \citet[323]{Calpestrati.2017} im Nachgang zu seiner sprachvergleichenden Korpusuntersuchung die Vermutung an, dass deutsche Muttersprachler:innen in der Schriftsprache native Intensivierer Fremd- bzw. Lehnwörtern gegenüber vorziehen. Ist dies bei meinem Korpus von Bestand, so ist anzunehmen, dass die Trefferzahl zuletzt genannter Ausdrücke grundsätzlich geringer ist als die von nativen Ausdrücken.

Die aufgeführten Punkte stellen zweifelsohne eine große Herausforderung für eine korpuslinguistische Untersuchung von Intensivierern dar, bei der die Korpusselektion umso bedeutender ist. Vor diesem Hintergrund stellt sich die durchaus berechtigte Frage, aus welchen Gründen ich trotz der offensichtlichen Mängel auf die zugrunde liegende Textform zurückgegriffen habe. Daher muss klargestellt werden, dass bislang keine umfangreichen, frei zugänglichen und linguistisch annotierten Korpora informeller Sprachvarietäten existieren, wie auch \citet[146]{Gutzmann.2019a} moniert. Infolge dieses methodischen Missstands greift er in seinen Untersuchungen häufig auf Google-Abfragen zurück; ein Vorgehen, das mir für meine Untersuchung weniger geeignet erscheint, als mit der mutmaßlich geringeren Ergiebigkeit des Korpus zu arbeiten. Daneben betont \citet[146]{Just.1967} unter Verweis auf die Wortschatzanalyse von \citet[]{Luck.1963}, dass die Journalist:innen des \textsc{Spiegels} selbst im Zuge der Berichterstattung oft umgangssprachlich formulierten, mit dem Ziel \glqq{}trockene Informationen plastisch und interessant\grqq{} zu machen, auch wenn das sprachliche Niveau dadurch schlimmstenfalls herabgesetzt werde. Dies lässt vermuten, dass die eigentlich konzeptionell schriftliche Interaktionsform des \textsc{Spiegels} dank anscheinend gezielt eingesetzter Umgangssprache ebenso in professionell redigierten Texten ausreichend Intensivierungen beinhaltet, deren diachronen Verläufe analysiert werden können. Sollten vereinzelte Intensivierer dennoch nur in geringer Anzahl auftreten, können durch andere jedoch idealerweise repräsentative Ergebnisse hervorgebracht werden. Ferner können positive Resultate in einem konzeptionell schriftlichen Korpus auch bei vergleichsweise wenigen Treffern nahelegen, dass die Ausdrücke in der gesprochenen Sprache fest etabliert sind und entscheidende Hinweise auf den gesprochenen Substandard geben.

Die vollständige, weitgehend ausgewogene Komposition des Zeitschriftenkorpus ist in Tabelle \ref{tab:SpZ} zusammengetragen.\footnote{Bei Prozentangaben handelt es sich hier und im weiteren Verlauf stets um gerundete Werte.} Diese Angaben entstammen restlos der Korpuspräsentation, die das \textsc{DeReKo} benutzerdefinierten Korpora bereitstellt (\citealt[]{IDS.2018}).

\begin{table}[p]
\begin{tabularx}{\textwidth}{lYYr}
\lsptoprule
Bezugsquelle & Texte & Wörter & Prozentualer Anteil \\
\midrule
\textsc{Der Spiegel} & 354\,773 & 238\,922\,315 & 49,17\,\% \\ 
\textsc{Die Zeit} & 366\,836 & 322\,353\,810 &   {50,83\,\%} \\
\midrule
In Summe & 721\,609 & 561\,276\,125 &   100\,\%\\
\lspbottomrule
\end{tabularx}
\caption[Vollständige Zusammensetzung des Zeitschriftenkorpus]{Vollständige Zusammensetzung des Zeitschriftenkorpus (in alphabetischer Reihenfolge).}
\label{tab:SpZ}
\end{table}

Anzumerken ist, dass die Ausgaben der \textsc{Zeit} von Januar 1994 bis November 1994 aufgrund von Nichtverfügbarkeit zur Gänze fehlen. In Anbetracht des verhältnismäßig kurzen Zeitraums, des leichten prozentualen Übergewichts des Teilkorpus sowie der anschließenden Normalisierung der Daten wird allerdings angenommen, dass das Fehlen dieser Ausgaben keine Auswirkungen auf die Resultate hat und demnach vernachlässigt werden kann. Obendrein gliedert sich die \textsc{Zeit} laut Korpuspräsentation wiederum in drei separate Bestandteile, wie Tabelle \ref{tab:Z} verdeutlicht. Diese ist nach dem Zeitraum sortiert, über den sich die jeweilige Bezugsquelle erstreckt.

\begin{table}[p]
\begin{tabularx}{\textwidth}{lYYrr}
\lsptoprule
{Bezugsquelle} & {Zeitraum} & {Texte} & {Wörter} & Prozentualer  {Anteil}\\
\midrule
\textsc{Die Zeit} & 1953--1998 & 235\,192 & 204\,657\,860 & 32,59\,\% \\
\textsc{Zeit - Online} & 1999--2014 & 107\,098 & 97\,427\,156 & 14,84\,\% \\
\textsc{Die Zeit} & 2015--2017 & 24\,546 & 20\,268\,794 & {3,40\,\%} \\
\midrule
In Summe & 1953--2017 & 366\,836 & 322\,353\,810 & 50,83\,\%\\
\lspbottomrule
\end{tabularx}
\caption[Zusammensetzung des Teilkorpus \textsc{Die Zeit}]{Zusammensetzung des Teilkorpus \textsc{Die Zeit} (angeordnet nach dem Zeitraum der Bezugsquelle).}
\label{tab:Z}
\end{table}

Das vollständige Belegkorpus enthält ganze Zeitschriftenausgaben und daher verschiedene Textsorten (bspw. Feuilletons, Dossiers, Interviews, Leserbriefe etc.), die entsprechend sprachnormativer Regeln redaktionell aufbereitet und professionell redigiert wurden. Deren Zusammensetzung ist in Tabelle \ref{tab:TeSo} dokumentiert. Diese Angaben wurden ebenfalls der Korpuspräsentation entnommen und sind gemäß abnehmender Textfrequenz aufgeführt.
%%%%%%

\begin{table}[p]
\begin{tabularx}{\textwidth}{lYY}
\lsptoprule
Textsorte & Texte & {Wörter}\\
\midrule
 \rowcolor[gray]{.90}
\textit{undefiniert} & 710\,814 & 552\,163\,847\\ 
\textit{Feuilleton} & 6\,338 & 	5\,266\,108 \\
 \rowcolor[gray]{.90}
\textit{Leserbrief} & 2\,183 & 	726\,799 \\
\textit{Dossier} & 1\,069 & 	1\,925\,430 \\
 \rowcolor[gray]{.90}
\textit{Interview} & 737 & 		861\,833 \\
\textit{Meldung} & 189 & 		159\,024 \\
 \rowcolor[gray]{.90}
\textit{Zitat} & 103 & 			15\,650 \\
\textit{Medienseite} & 66 & 	60\,991 \\
 \rowcolor[gray]{.90}
\textit{Kommentar} & 60 & 		37\,846 \\
\textit{Essay} & 32 & 			57\,004 \\
 \rowcolor[gray]{.90}
\textit{Bildbeischrift} & 16 & 	188 \\
\textit{Beilage} & 2 & 			1\,405 \\
\midrule
In Summe & 721\,609 & 561\,276\,125\\
\lspbottomrule
\end{tabularx}
\caption[Textsortenkomposition im Zeitschriftenkorpus]{Zusammensetzung der Textsorten im Zeitschriftenkorpus (angeordnet nach abnehmender Frequenz).}
\label{tab:TeSo}
\end{table}

\noindent Wie Tabelle \ref{tab:TeSo} zu entnehmen ist, ist der Großteil des Zeitschriftenkorpus nicht definiert. Aus Gründen der Vollständigkeit sei erwähnt, dass somit auch andere als die hier aufgeführten Textsorten enthalten sein können. Genauso ist es möglich, dass sich welche der darunter dargelegten Textsorten (z.\,B. Interview, Meldung oder Zitat) in dieser Gruppe wiederfinden, ohne dass sie in den Metadaten entsprechend zugeordnet sind. Ob diese Mutmaßung der Realität entspricht, lässt sich aufgrund fehlender Klassifikation jedoch nicht abschätzen.

Insgesamt setzt sich das virtuelle Arbeitskorpus nach Angaben der Korpuspräsentation aus 68 Jahrgängen, bestehend aus 133 Subkorpora mit 721\,609 Texten und ca. 562 Millionen Wörtern, zusammen. Die Korpuszusammensetzung je Dekade ist in chronologischer Reihenfolge in Tabelle \ref{tab:ZuSa} aufgelistet.

%
\begin{table}
\begin{tabularx}{\textwidth}{lYY}
\lsptoprule
Jahrzehnt & Texte & Wörter \\
\midrule
 \rowcolor[gray]{.90}
1950--1959 & 53\,424 & 34\,475\,290\\ 
1960--1969 & 91\,899 & 66\,438\,193 \\
 \rowcolor[gray]{.90}
1970--1979 & 104\,608 & 78\,524\,852 \\
1980--1989 & 109\,608 & 98\,801\,677 \\
 \rowcolor[gray]{.90}
1990--1999 & 121\,545 & 93\,748\,502 \\
2000--2009 & 143\,253 & 107\,160\,066\\
 \rowcolor[gray]{.90}
2010--2019 & 97\,272 & {82\,127\,545} \\
\midrule
In Summe & 721\,609 & 561\,276\,125\\
\lspbottomrule
\end{tabularx}
\caption[Quantitative Zusammensetzung des Zeitschriftenkorpus]{Quantitative Zusammensetzung des Zeitschriftenkorpus in Jahrzehnten (in chronologischer Reihenfolge).}
\label{tab:ZuSa}
\end{table}

Der anschließende Ausschnitt aus der \textsc{Zeit} exemplifiziert einen zufällig ausgewählten Textausschnitt eines Berichts über Steigerungspartikeln im Deutschen und dient dazu, einen kurzen Einblick in die Schreibweise einer Textsorte des Untersuchungskorpus zu gewähren:\footnote{Neben der Illustration eines prototypischen Zeitschriftenberichts ist der ausgewählte Text auch aus zwei weiteren Gründen interessant: Zum einen spielt er auf den in diesem Buch in H1 untersuchten Abnutzungseffekt an, dem Intensivierer dem Anschein nach unterliegen, der dafür sorge, dass sie mit der Zeit die angestrebte Ausdrucksstärke verlieren \citep[vgl.][2]{Hauschild.1899}. Zum anderen findet sich im genannten Beispiel \glqq{}Er ist total der Spinner\grqq{} die in Abschnitt \ref{Intensivierertypen} besprochene Möglichkeit zur DP-externen Realisierung des expressiven Intensivierers, die u.\,a. \citet[]{Schlieben.1995} beschreibt.} 
 
 \begin{quote}
 Man muß zugeben, daß unsere Sprache bis zu dieser Entdeckung unter einem bedenklichen Mangel an glaubwürdigen Steigerungsformeln litt. Lange Zeit beliebte Kombinationen mit dem Wort \glqq{}total\grqq{}, wie sie in der Verbindung \glqq{}Er ist total der Spinner\grqq{} zum Ausdruck kommen, sind ja selbst unter Heranwachsenden im Abklingen begriffen. Das bedingt verwendbare \glqq{}absolut\grqq{} ist für Interviews nach gewonnenen Matches reserviert und gehört inzwischen Steffi Graf allein. Selbst das lockere \glqq{}super\grqq{} hat die alte Strahlkraft verloren, seit es in den Vokabelschatz von Entzücken heuchelnden Boutique-Besitzerinnen und sich jugendlich gebenden Managern abgesunken ist.   \hfill (\mbox{\textsc{Die Zeit} -- Z89/MAR.00044 -- 03.03.1989)}
\end{quote}

\subsection{Belegsammlung}
Basierend auf Alltagsbeobachtungen wurden zwischen Juni 2018 und Januar 2019 stichprobenartig ausgewählte Ausdrücke (n = 25) einzeln über die Suchmaske der Webapplikation COSMAS II abgerufen. Dabei handelte es sich überwiegend um Hörbelege aus der gesprochenen Alltagssprache; die daraus extrahierten Intensivierer sind auch in der einschlägigen Literatur dokumentiert (vgl. bspw. \citealt[]{Biedermann.1969} und \citealt[]{Os.1989}). Die Menge der zufällig definierten Stichprobe wurde zu Beginn der Untersuchung bewusst offengehalten, um der beachtlichen Variation der deutschen Intensivierer wenigstens ansatzweise gerecht zu werden, ohne Anspruch auf Vollständigkeit zu erheben. Dabei wurden andere im Deutschen potenziell relevante Kandidaten nicht explizit ausgeschlossen, sie traten im beobachteten Datenmaterial jedoch schlicht nicht auf.\footnote{Dies betraf u.\,a. den multifunktionalen Intensivierer \textit{so}, nach dessen Fehlen sich ein:e anonyme:r Gutachter:in erkundigt hatte. Dieser Ausdruck wurde nicht bewusst von der Untersuchung ausgeschlossen, sondern fand sich in keinem der zuvor erhobenen Hörbelege -- gleichwohl er neueren Studien zufolge (bspw. \citealt[]{Stratton.2020b}; \citealt[]{Scheffler.2023}) auch in sozialen Medien zu den häufigsten Intensivierern zu gehören scheint.} Die maschinelle Abfrage beinhaltete die folgenden expressiven Intensivierer (n = 24):\footnote{Die Angabe "'n ='' bezieht sich hier und im weiteren Verlauf auf die Anzahl der untersuchten Ausdrücke pro Intensiviererklasse und nicht auf deren Vorkommen im Korpus.} \textit{arg}, \textit{arsch}, \textit{derb}, \textit{end}, \textit{fürchterlich}, \textit{furchtbar}, \textit{hammer}, \textit{krass}, \textit{mega}, \textit{ober}, \textit{sau}, \textit{scheiß}, \textit{scheiße}, \textit{schrecklich}, \textit{schweine}, \textit{super}, \textit{übel}, \textit{übelst}, \textit{übertrieben}, \textit{ultra}, \textit{unheimlich}, \textit{verdammt}, \textit{voll} und \textit{wahnsinnig}. Diese wurden einerseits aufgrund ihrer subjektiv wahrnehmbaren Aktualität ausgewählt (z.\,B. \textit{krass} und \textit{mega}), andererseits aber auch deswegen, weil sie als usualisierte und daher idiomatisierte Intensivierer des Deutschen gelten, die in der Umgangssprache häufig vorkommen (z.\,B. \textit{sau} und \textit{voll}). Überdies weisen einige der Intensivierer eine offenkundige steigernde Semantik auf (z.\,B. \textit{ober} oder \textit{über}). Somit umfasst das Untersuchungsmaterial der semantischen Herkunft nach sowohl native, transparente Ausdrücke mit klar ersichtlichem Ursprung als auch opake Ausdrücke, deren Quelllexem nicht (mehr) wahrnehmbar ist. Neben den eher jugendsprachlichen Intensivierern \textit{end}, \textit{hammer} und \textit{übel(st)} wurden verschiedene Ausdrücke mit klar negativer Wortherkunft in die Untersuchung miteingebunden, die u.\,a. \citet[165ff.]{Biedermann.1969} und \citet[]{Hentschel.1998b} besprechen (z.\,B. \textit{furchtbar} und \textit{schrecklich}). Insgesamt ist also ein Großteil der von \citet[]{Claudi.2006} angeführten semantischen Quellbereiche mit mindestens einem Stellvertreter (z.\,B. \textit{wahnsinnig} für die Gruppe \textit{Mentale Erkrankungen} oder \textit{verdammt} für die Gruppe \textit{Religion}) abgedeckt.

Dadurch ergibt sich eine vielseitige Auswahl, die einen guten, wenn auch angesichts der schieren Menge an Ausdrücken selbstverständlich keineswegs erschöpfenden Überblick über die im Deutschen verwendeten expressiven Intensivierer bietet, mit dessen Hilfe die Korrektheit der in der theoretischen Literatur vertretenen Hypothesen H1 und H2 empirisch überprüft werden kann.

Als Ergänzung zu den expressiven Intensivierern wurde \textit{sehr} als klassischer Vertreter der stabilen deskriptiven Intensiviererklasse miteinbezogen (n = 1), der bei der Auswertung als Vergleichsbasis fungiert, an der die diachronen Entwicklungsverläufe der expressiven Intensivierer abgeglichen werden können.\footnote{Ein:e anonyme:r Gutachter:in beanstandet, dass es durchaus problematisch sein könne, aus einem einzelnen Ausdruck Rückschlüsse auf die gesamte deskriptive Intensiviererklasse zu ziehen. Zudem weist die Person darauf hin, dass die Entwicklung von deskriptiven Ausdrücken wie \textit{ganz}, \textit{so} oder \textit{ziemlich} ebenso von Interesse sein kann. Da in dieser Untersuchung jedoch die expressive Intensiviererklasse im Fokus steht und nicht die allgemeine Entwicklung deskriptiver Intensivierer im Korpus, halte ich diesen Einwand für weniger entscheidend – zumal \textit{so} multifunktional ist und \textit{ganz} und \textit{ziemlich} eine andere skalare Funktion aufweisen als \textit{sehr} (vgl. Abschnitt \ref{Funktionsklasse})}. Diesbezüglich wird erwartet, dass \textit{sehr} infolge seiner rein skalaren und neutralen Bedeutung im Zeitschriftenkorpus erheblich frequenter ist als die expressiven Intensivierer. So berichten bspw. \citet[266]{Ito.2003}, dass das englische Pendant \textit{very} in ihrer Untersuchung der mit Abstand am häufigsten gebrauchte Intensivierer im Korpus war. Nach intuitiver Einschätzung trifft auf den Gebrauch von \textit{sehr} in den hier analysierten Zeitschriften Ähnliches zu.\footnote{An dieser Stelle sei erwähnt, dass sämtliche Belege mit allen syntaktischen Verwendungsarten (adjektivisch, adverbial und verbal) ungefiltert in die Untersuchung einbezogen wurden, sofern es sich um Intensivausdrücke handelte. Auf eine quantitative Auswertung ihrer jeweiligen Vorkommenshäufigkeit im Korpus sowie der modifizierten Köpfe wurde verzichtet; eine nachträgliche Erhebung dieser Daten ist im Rahmen der vorliegenden Untersuchung nicht mehr umsetzbar. Zur weiteren Erhöhung der Validität sollten diese Aspekte in künftigen Studien bedacht und gezielt berücksichtigt werden.}

Tabelle \ref{tab:Int} bildet die mittels des Korpus untersuchten Intensivierer inklusive ihrer semantischen und lexikalischen Merkmale ab. Diese sind nach semantischem Herkunftsbereich sortiert. Abgesehen von der Klassifikation ihres Bedeutungstyps ist angegeben, ob die Wortherkunft des jeweiligen Ausdrucks opak oder transparent ist, was im Hinblick auf die spätere Auswertung von Interesse sein kann. Außerdem ist aufgeführt, in welcher Sprache bzw. Zeitperiode des Deutschen der etymologische Ursprung zu verorten ist. Diese Einordnung erfolgt beinahe ausschließlich auf Basis von \citet[]{Pfeifer.1995}; lediglich bei \textit{übertrieben} wurde aufgrund von Nichtverfügbarkeit zur Informationsgewinnung auf das \citet[607f.]{Grimm.1956} zurückgegriffen. Im Sinne der Einordnung nach \citet[354ff.]{Claudi.2006} wird ergänzend die semantische Abstammung der Ausdrücke (vgl. Abschnitt \ref{Quellensemantisch}) sowie die mit ihnen verbundene konnotative Markiertheit berücksichtigt. Letztere ist in der Tabelle jeweils mit `Positiv', `Neutral' oder `Negativ' gekennzeichnet. Diese Klassifikation wird aus der Annahme heraus vorgenommen, dass die evaluative Polarität eines Ausdrucks gegebenenfalls Einfluss auf die Frequenz in schriftsprachlichen Kontexten hat. So ist davon auszugehen, dass der Gebrauch eines transparenten Intensivierers wie \textit{arsch} oder \textit{scheiß} infolge negativer Semantik und ebensolcher Konnotation im Rahmen einer eher konservativen Zeitschrift grundsätzlich weniger akzeptabel ist als der Gebrauch eines neutralen Intensivierers wie \textit{sehr}, eines weniger offensiven Intensivierers wie \textit{arg}, eines opaken Intensivierers wie \textit{krass} oder eines positiv konnotierten Intensivierers wie \textit{super}. Aus diesem Grund lassen sich derartige offensive und stark negativ konnotierte Ausdrücke im Zeitschriftenkorpus mutmaßlich seltener nachweisen als stilistisch neutralere, textsortenadäquatere Ausdrücke. In diesem Sinne ist also vorstellbar, dass sich die semantische Herkunft sowie die häufig damit einhergehende Konnotation eines Intensivierers auf dessen Verwendung auswirken und diese Faktoren unter Umständen auch dazu in der Lage sind, dessen Häufigkeit im Korpus zu erklären.

\begin{sidewaystable}
\small
\begin{tabular}{llllll}
\lsptoprule
Intensivierer & Klassifikation & Semantik & Datierung/Herkunft & Herkunftsbereich & Konnotation\\
\midrule
\rowcolor[gray]{.90}
\textit{sehr} & deskriptiv & opak & Althochdeutsch & Masse \& Kraft & Neutral\\ 				
\textit{arg} & expressiv & transparent & Althochdeutsch  & Leid \& Furchteinflößendes & Negativ\\
\rowcolor[gray]{.90}
\textit{fürchterlich} & expressiv & transparent & Althochdeutsch & Leid \& Furchteinflößendes & Negativ\\		
\textit{furchtbar} & expressiv & transparent & Mittelhochdeutsch & Leid \& Furchteinflößendes & Negativ\\ 
\rowcolor[gray]{.90}
\textit{schrecklich} & expressiv & transparent & Frühneuhochdeutsch & Leid \& Furchteinflößendes & Negativ\\ 
\textit{übel} & expressiv & transparent & Althochdeutsch & Leid \& Furchteinflößendes & Negativ\\	
\rowcolor[gray]{.90}			
\textit{übelst} & expressiv & transparent & Althochdeutsch & Leid \& Furchteinflößendes & Negativ\\ 
\textit{unheimlich} & expressiv & transparent & Mittelhochdeutsch & Leid \& Furchteinflößendes & Negativ\\
\rowcolor[gray]{.90}	
\textit{arsch} & expressiv & transparent & Althochdeutsch & Körperliches & Negativ\\
\textit{scheiß} & expressiv & transparent & Mittelhochdeutsch & Körperausscheidungen/Fäkalien & Negativ\\	
\rowcolor[gray]{.90}
\textit{scheiße} & expressiv & transparent & Mittelhochdeutsch & Körperausscheidungen/Fäkalien & Negativ\\ 
\textit{wahnsinnig} & expressiv & transparent & Frühneuhochdeutsch & Mentale Erkrankungen & Negativ \\ 
\rowcolor[gray]{.90}
\textit{verdammt} & expressiv & transparent & Althochdeutsch & Religion & Negativ\\	
\textit{sau} & expressiv & transparent & Althochdeutsch & Animalisches & Negativ\\ 	
\rowcolor[gray]{.90}			
\textit{schweine} & expressiv & transparent & Althochdeutsch & Animalisches & Neutral\\	
\textit{derb} & expressiv & transparent & Althochdeutsch & Masse \& Kraft & Negativ\\
\rowcolor[gray]{.90}
\textit{hammer} & expressiv & transparent & Althochdeutsch & Masse \& Kraft & Neutral\\		
\textit{krass} & expressiv & opak & Lateinisch & Masse \& Kraft & Neutral\\			
\rowcolor[gray]{.90}
\textit{mega} & expressiv & opak & Griechisch & Masse \& Kraft & Positiv\\
\textit{voll} & expressiv & transparent & Althochdeutsch & Masse \& Kraft & Neutral\\	
\rowcolor[gray]{.90}	
\textit{end} & expressiv & transparent & Althochdeutsch & Endpunkt/Schwellenwert & Neutral \\			
\textit{ober} & expressiv & transparent & Althochdeutsch & Endpunkt/Schwellenwert & Neutral\\ 
\rowcolor[gray]{.90}  		 	
\textit{super} & expressiv & opak & Lateinisch & Endpunkt/Schwellenwert & Positiv\\ 					
\textit{übertrieben} & expressiv & transparent & Mittelhochdeutsch & Endpunkt/Schwellenwert & Negativ\\  
\rowcolor[gray]{.90}	
\textit{ultra} & expressiv & opak & Lateinisch & Endpunkt/Schwellenwert & Neutral\\
\lspbottomrule
\end{tabular}
\caption[Gesamtübersicht der korpuslinguistisch untersuchten Intensivierer]{Gesamtübersicht der korpuslinguistisch untersuchten Intensivierer (n = 25) und deren semantischen sowie lexikalischen Merkmalen (angeordnet nach semantischem Herkunftsbereich).}
\label{tab:Int}

\end{sidewaystable}

%%%%%%%%%%%%%%%%%%%%%%%%%%%

\subsection{Belegfilterung}
\largerpage
Die Filterung der Korpusbelege wurde inkrementell vorgenommen. Im ersten ergebnisoffenen Schritt ging es zunächst um die maschinelle Kompletterfassung derjenigen Sätze, die die zu untersuchenden Ausdrücke in irgendeiner Form -- ganz gleich, ob als Adjektiv, Adverb oder Nomen -- beinhalteten. Dies war dem ungünstigen Umstand geschuldet, dass die Annotation des Untersuchungskorpus strukturell und funktionsmäßig nicht auf die Fragestellung abgestimmt war, sodass eine gezielte Suche nach Intensivierern nicht möglich war. Durch die fehlertolerante Abfrage enthielten die exportierten Listen gleichzeitig eine Unmenge an nicht-intendierten falsch positiven Treffern -- sogenannte \glqq{}Fehler 1. Art\grqq{} bzw. \glqq{}$\alpha$-Fehler\grqq{} (\citealt[62]{Perkuhn.2012}) -- , die durch die Korpusabfrage zwar gefunden wurden, die für die zu klärenden Hypothesen jedoch nicht von Interesse waren und die es daher zu beseitigen galt. Aus diesem Grund wurden im zweiten Schritt alle Vorkommen der extrahierten Ausdrücke mithilfe der computerlinguistischen Annotationssoftware \textit{TreeTagger} (\citealt[]{Schmid.1994,Schmid.1995}) automatisch morphosyntaktisch getaggt. Dieser Part-of-speech-Tagger ist für mehrere Sprachen (bspw. Englisch oder Französisch) erhältlich und annotiert etwa 8\,000 Tokens pro Sekunde mit einer Erfolgsrate von ca. 97,53\,\% \citep[vgl.][22f.]{Schmid.1995}. Die verbliebenen Treffer wurden im letzten Schritt manuell in Bezug auf Intensivierungen hin bereinigt. Infolge dieser schrittweisen und gründlichen Filterung konnten nicht-verwertbare Fehlerquellen schon von Beginn der Materialauswertung an ausgeschlossen werden. Zu solchen gehörten u.\,a. Nomen, deren Basen nicht ohne Bedeutungsänderung frei stehen können (z.\,B. \textit{Megalith} oder \textit{Megaphon}), Nominalkomposita (z.\,B. mit \textit{Super-} wie in \textit{Superbenzin} oder \textit{Supermann}), Markenbezeichnungen (z.\,B. \textit{Superdry} oder \textit{Super-Dickmanns}) sowie Toponyme und Choronyme (z.\,B. \textit{Hammerfest} oder \textit{Sauerland}). Im Falle von \textit{voll(-)} wurden absolute, nicht-expressive Verwendungen wie in \textit{voll ausgelastet}, \textit{vollbeschäftigt} oder \textit{voll funktionsfähig} aussortiert. Daneben wurden unerwünschte Doppelbelege und rein qualitative Verwendungen ausgeklammert, bspw. Kombinationen mit Partizipien wie in \textit{arg gebeutelt} oder \textit{arg zugesetzt}. Darüber hinaus wurden alle anderen adjektivischen Verwendungen der Ausdrücke (z.\,B. \textit{voll geistiger Unruhe}) bzw. solche, in denen die Intensivierer als Wortbestandteil vorkommen und die bloß aufgrund der ergebnisoffenen Suche in der Trefferliste gelandet sind, ausgesiebt (z.\,B. \textit{sauber} oder \textit{vollschlank}).

Um das Annotationsvorgehen transparent zu machen, werden im Anschluss zufällig ausgewählte Sätze vorgestellt, die entweder ausgeklammert oder in die Auswertung miteingebunden wurden. Begonnen wird in \REF{ausschluss} mit falsch positiven Treffern, die für die spätere Analyse nicht infrage kamen.

\largerpage
\ea \label{ausschluss}
\ea Auch der Preis ist superlativisch. \\   	\label{preis}
\textsc{Die Zeit} -- Z73/JUL.00285 -- 20.07.1973
\ex Pädagogendeutsch ist schrecklich, erst recht bei weltanschaulich gebundenen Pädagogen.  \\   \label{päd}
\textsc{Die Zeit} -- Z66/AUG.00057 -- 05.08.1966
\ex Wenn man es kraß ausdrücken will -- und da der Verfasser selber zu dieser Gruppe und somit zu den Kritisierten gehört, ist ihm eine krasse Formulierung vielleicht erlaubt --, haben sie ihre `Erstgeburt für ein Linsengericht', haben sie ihr intellektuelles Niveau für ein gutes Gehalt verkauft. \\  \label{krasse}
\textsc{Der Spiegel} -- S65/APR.00259 -- 21.04.1965
\z
\z
Während das Adjektiv \textit{superlativistisch} in Satz \REF{preis} nicht in \textit{super-} und \textit{-lativistisch} aufgespalten werden kann, \textit{super-} in diesem Fall folglich keine modifizierende Funktion hat, handelt es sich bei \textit{schrecklich} in Satz \REF{päd} um ein prädikatives Adjektiv, das ebenso aussortiert wurde. In Satz \REF{krasse} findet sich mit \textit{kraß} zum einen eine modale adverbiale Bestimmung und zum anderen mit \textit{krasse} ein attributives Adjektiv. Auch solche Belege wurden im Zuge der Filterung ausgeschlossen. Kurz gefasst wurde also all das Material außer Acht gelassen, bei dem es sich nicht um Intensivausdrücke in Form von Intensivierer-Adjektiv-Kombinationen handelte. In die Beleganalyse einbezogen wurden hingegen alle attributiven \REF{ziga}, prädikativen \REF{sitz} und adverbialen Vorkommen \REF{langhaar} von Intensivausdrücken, aber auch Ausrufe wie in \REF{ausruf}, (rhetorische) Fragen \REF{komfort} sowie Stichwörter \REF{sauer}, wie an beispielhaft ausgewählten Belegen von \textit{super} vorgeführt wird:

\ea \label{subba}
\ea Eine superlange Zigarette mit altem Image soll den Markt erobern.  \\  \label{ziga}
\textsc{Die Zeit} -- Z68/SEP.00337 -- 20.09.1968
\ex Die Sitzgruppe ist superwuchtig und von amerikanischer Bequemlichkeit, ein weiß gestrichenes schmiedeeisernes Gitter und eine großflächige Bildertapete teilen die Eßgarnitur vom übrigen Wohnraum. \\  \label{sitz}
\textsc{Der Spiegel} -- S65/AUG.00099 -- 11.08.1965
\ex Da reitet ein langhaariger Dschihadist mit über die Schulter gehängtem Sturmgewehr auf einem Araberhengst -- sieht supercool aus.\\  \label{langhaar}
\textsc{Die Zeit} -- Z15/JAN.00438 -- 20.01.2015
\ex Die Botschaft ist immer die gleiche: supereinfach, superlecker!\\  \label{ausruf}
\textsc{Der Spiegel} -- S15/DEZ.00067 -- 05.12.2015
\ex Superkomfortabel?\\  \label{komfort}
\textsc{Die Zeit} -- Z17/SEP.00322 -- 21.09.2017
\ex Gemütszustand: supersauer. \\  \label{sauer}
\textsc{Die Zeit} -- Z16/SEP.00803 -- 29.09.2016
\z
\z
Wie in Beispiel \REF{subba} zu sehen ist, modifiziert \textit{super} in allen Sätzen das adjazente graduierbare Adjektiv. Demnach handelt es sich um einen Intensivierer, einen positiven Treffer also, der für die Analyse von Interesse war.

Die graphematische Repräsentation der Intensivierer-Adjektiv-Kombination spielte für die Annotation keine Rolle: Morphologisch inkorporierte (\textit{megagut}) wie auch isolierte (\textit{mega gut}) oder mittels Bindestrich miteinander verbundene Varianten (\textit{mega-gut}) sind in die Auswertung eingeflossen, sofern die Partikel unstrittig mit einer Gradbedeutung einherging. Vor diesem Hintergrund ist hervorzuheben, dass in Bezug auf die Schreibweise von Intensivausdrücken gemeinhin kein Konsens besteht. Empfiehlt das vom Rat für deutsche Rechtschreibung herausgegebene amtliche Regelwerk\footnote{\citet[38]{Rfdr.2018}:
  \url{https://www.rechtschreibrat.com/DOX/rfdr_Regeln_2016_redigiert_2018.pdf} [letzter Aufruf: 25.06.2021].
} (§\,36,\,Abs.\,1.5) sowohl im Falle von intensitätsverstärkenden wie auch -mindernden Ausdrücken grundsätzlich die Zusammenschreibung, trifft man in der theoretischen Literatur auf eine große Variabilität. So weist bspw. \citet[203]{Gualberto.2011} darauf hin, dass Gradadjektive wie \textit{hoch qualifiziert} oder \textit{tief erschüttert} nach den Regeln der neuen deutschen Rechtschreibung ebenso als Juxtaposition getrennt vom modifizierten Adjektiv geschrieben werden können. Dies sei vor allem dann möglich, wenn sie zusammen mit adjektivisch verwendeten Partizipien auftreten, wohingegen bei rein intensivierenden Verwendungsweisen Zusammenschreibung zu empfehlen sei, z.\,B. \textit{hochintelligent} oder \textit{tieftraurig} \citep[vgl.][203]{Gualberto.2011}. Dagegen beobachtet \citet[]{Renz.2020} die Getrenntschreibung in erster Linie bei Intensivierern mit eigenem Wortstatus. So treten in seinem Korpus vorzugsweise diejenigen Intensivierer abgetrennt vom Adjektiv auf, die entweder sowieso eigenständige lexikalische Einheiten darstellen, wie \textit{furchtbar} oder \textit{schrecklich}, oder infolge hoher Produktivität als ebensolche interpretiert werden können, z.\,B. \textit{mega} oder \textit{super}. Ob ein Ausdruck bloß graduierend im Sinne von skalierend ist oder er aufgrund seines Wortstatus selbständig auftreten kann, macht aus meiner Sicht für die Zielsetzung der in diesem Buch untersuchten Adjektiv-Intensivierung aus mehreren Gründen keinen Unterschied:\footnote{Anders
  sieht dies in den Fällen aus, in denen expressive Ausdrücke wie \textit{mega} oder \textit{super} in adjektivischer Verwendung adjazente Nomen (z.\,B. \textit{mega Party} oder \textit{Supermarkt}) näher beschreiben. Wie \citet[]{Ehrmantraut.2022} empirisch begründet, macht es bei Nomen durchaus einen Unterschied, ob ein solches Morphem inkorporiert oder infolge von morphologischem und semantischem Wandel selbständig auftritt. So verfügen letztere, d.\,h. vormals gebundene Morpheme, über eine neue expressive Bedeutung, die zugleich eine semantische Verengung nach sich ziehe: \textit{Die Party ist mega.} vs. \textit{\#Der Markt ist super}.
} Einerseits wird die durch das Adjektiv ausgesagte Eigenschaft bei allen dieser Schreibvarianten auf der Skala höher positioniert, was im Vergleich zum Positiv einen verstärkten Intensitätsgrad nach sich zieht. Andererseits steht hier primär die Frequenz der Ausdrücke in intensivierender Funktion im Korpus im Fokus. Und Intensivierung liegt in allen Fällen vor -- ungeachtet der graphischen Schreibweise. Außerdem gehen zahlreiche Autor:innen bzw. Journalist:innen meinem Empfinden nach diesbezüglich ohnehin mehr intuitiv als regelgeleitet vor, was in Abhängigkeit des Intensivierers zu unterschiedlichen Schreibweisen führen kann. Aus diesen Gründen habe ich alle drei Varianten (d.\,h. gebunden, mit Spatium oder Durchkopplungsbindestrich) in die Analyse miteingebunden. Um herauszufinden, ob die Partikel eine graduierende Bedeutung hat, wurde der in Abschnitt \ref{Intensivierertypen} eingeführte Paraphrasierungstest mit \textit{sehr} angewandt, der u.\,a. in der von \citet[]{Klara.2009} durchgeführten Korpusstudie zum Einsatz kam. Dabei wird geprüft, ob sich der jeweilige expressive Intensivierer durch das deskriptive \textit{sehr} ersetzen lässt \citep[vgl.][191]{Wellmann.1978}; ist dies der Fall, so wird der Belegsatz miteinbezogen, andernfalls ausgeklammert.\footnote{Für eine ausführliche Diskussion der Vor- und Nachteile der \textit{sehr}-Paraphrase sei auf \citet[]{Klara.2009} verwiesen.}

Da in einigen Fällen trotz manueller Annotation und Kontextbegutachtung keine eindeutige Zuordnung und folglich keine Auflösung der bestehenden Funktionsambiguitäten erfolgen konnte, wurden einige der expressiven Intensivierer aus methodischen Gründen von der Analyse ausgeschlossen. Dies betraf zum einen die funktionsambigen Ausdrücke \textit{übel}, \textit{übelst}, \textit{übertrieben} und \textit{ultra} (n = 4), die in vielen Fällen mehrere Lesarten zulassen und hinsichtlich der intensivierenden Funktion nicht zweifelsfrei trennscharf kategorisierbar waren.\footnote{Ähnliches berichtet \citet[557]{Athanasiadou.2007} in Bezug auf den englischen Intensivierer \textit{quite}, der ebenso nur vage interpretierbar sei und daher gesondert betrachtet werde müsse.} Zum anderen wurden \textit{arsch}, \textit{end}, \textit{scheiße} und \textit{schweine} (n = 4) ausgeklammert, die im Korpus zu niedrigfrequent waren und infolgedessen auf wenig erfolgversprechende Ergebnisse hoffen ließen. Aufgrund ihres seltenen Vorkommens im Korpus lassen sich diese Ausdrücke vielmehr als Zufallsfunde denn als systematisch verwendete Intensivierer klassifizieren. Hervorzuheben ist, dass es sich bei ihnen keineswegs um neue Intensivierer handelt, die vermeintlich in der Frequenz zunehmen, sondern möglicherweise um zufällige idiosynkratische okkasionelle Bildungen, die (zumindest in diesem Korpus) nicht weiter von Bedeutung sind. Um der Zielsetzung entsprechend eine repräsentative Stichprobe zu erhalten, wurden also acht Intensivierer vor der Analyse ausgesiebt, die in Summe 4\,105 absolute Treffer und somit ca. 21,15\,\% der Gesamtbelege ausmachten. Nach dem Ausschluss dieser Ausdrücke blieben 16 expressive Intensivierer übrig, die im Fokus der nachfolgenden Ausführungen stehen: \textit{arg}, \textit{derb}, \textit{furchtbar}, \textit{fürchterlich}, \textit{hammer}, \textit{krass}, \textit{mega}, \textit{ober}, \textit{sau}, \textit{scheiß}, \textit{schrecklich}, \textit{super}, \textit{unheimlich}, \textit{verdammt}, \textit{voll} und \textit{wahnsinnig}. Tabelle \ref{tab:HäuVer} bietet eine vollständige quantitative Übersicht aller aus dem Korpus extrahierten Intensivierer.

\begin{table}
\begin{tabularx}{\textwidth}{XlYY}
\lsptoprule
{Intensivierer} & {Intensiviererklasse} & \multicolumn{2}{c}{\textbf{Korpusfrequenz}}\\
 & & absolut & relativ\\
\midrule
\rowcolor[gray]{0.90}
\textit{sehr} & deskriptiv & 165\,161 & 20\,278,12\\ 
\textit{super} & expressiv & 3\,476 & 391,26 \\
\rowcolor[gray]{0.90}
\textit{ultra} & expressiv & 3\,243 & 374,10 \\
\textit{verdammt} & expressiv & 2\,020 & 228,42 \\
\rowcolor[gray]{0.90}
\textit{schrecklich} & expressiv & 1\,600 & 186,26 \\
\textit{arg} & expressiv & 1\,594 & 185,50 \\
\rowcolor[gray]{0.90}
\textit{furchtbar} & expressiv & 1\,557 & 177,11 \\
\textit{unheimlich} & expressiv & 1\,427 & 161,89 \\
\rowcolor[gray]{0.90}
\textit{wahnsinnig} & expressiv & 1\,351 & 137,58 \\
\textit{übertrieben} & expressiv & 732 & 91,50 \\
\rowcolor[gray]{0.90}
\textit{sau} & expressiv & 523 & 60,41 \\
\textit{fürchterlich} & expressiv & 397 & 43,48 \\
\rowcolor[gray]{0.90}
\textit{scheiß} & expressiv & 403 & 41,97 \\
\textit{voll} & expressiv & 357 & 36,65 \\
\rowcolor[gray]{0.90}
\textit{derb} & expressiv & 149 & 22,06 \\
\textit{ober} & expressiv & 185 & 20,53 \\
\rowcolor[gray]{0.90}
\textit{mega} & expressiv & 115 & 11,43 \\
\textit{hammer} & expressiv & 78 & 7,70 \\
\rowcolor[gray]{0.90}
\textit{krass} & expressiv & 73 & 7,11 \\
\textit{übel} & expressiv & 53 & 6,14 \\
\rowcolor[gray]{0.90}
\textit{schweine} & expressiv & 29 & 2,91 \\
\textit{arsch} & expressiv & 27 & 2,68 \\
\rowcolor[gray]{0.90}
\textit{übelst} & expressiv & 16 & 1,57 \\
\textit{end} & expressiv & 3 & 0,29 \\
\rowcolor[gray]{0.90}
\textit{scheiße} & expressiv & 2 & 0,19 \\
\midrule
In Summe & & 184\,570 & 22\,476,86\\
\lspbottomrule
\end{tabularx}
\caption[Absolute und relative Häufigkeiten der korpuslinguistisch untersuchten Intensivierer]{Absolute und relative Häufigkeiten der 25 untersuchten Intensivierer im Zeitraum von Januar 1950 bis Dezember 2017 bei einer Korpusgröße von ca. 562 Millionen Textwörtern (in abnehmender relativer Frequenz).}
\label{tab:HäuVer}
\end{table}

Nachdem die exportierten, tabellarisch gelisteten Belege inklusive Kontextsatz, Sigle, Datum und Quellenangabe manuell auf Intensivierungen hin überprüft und die oben angemerkten falsch positiven Treffer beseitigt wurden, konnten sie für die weitere Datenauswertung genutzt werden. Die sorgfältige Belegfilterung und -aufbereitung gewährte es, dass ausschließlich diejenigen Treffer in die tiefergehende Analyse eingeflossen sind, die tatsächlich relevant für die Beantwortung der eingangs formulierten Forschungsfragen waren. Die Trefferliste eines jeden Intensivierers entspricht seiner absoluten Gesamtzahl. Diese absolute Häufigkeit bildet das reine Vorkommen des Ausdrucks im Korpus ab, die verständlicherweise von der jährlichen Korpusgröße bzw. vom Textumfang der Zeitschriften abhängig ist. Da dieser Umfang über den Zeitraum von 68 Jahren allerdings stark schwankt, wurden die absoluten Werte mit Blick auf die Zielsetzung der Untersuchung auf eine Million Wörter (englisch: \textit{instances per million words}, hier und im Folgenden abgekürzt mit \textit{IpMW}) normalisiert. Durch diese Skalierung konnte das anfangs bestehende Ungleichgewicht ausbalanciert und diachrone Vergleichbarkeit sichergestellt werden. Dieser relative Frequenzwert gibt an, wie viele Instanzen eines Ausdrucks in einer Million Textwörter des Untersuchungskorpus enthalten sind. Im Anschluss daran wurden die longitudinale Gesamtentwicklung der Intensivierer unter Zuhilfenahme von Histogrammen grafisch sichtbar gemacht und darauf aufbauend deren quantitativen Entwicklungsverläufe miteinander verglichen.

Die Ergebnisse der Korpusuntersuchung, die sich, sofern nicht anders angegeben, auf die normalisierten, d.\,h. relativen Frequenzen in Form von IpMW stützen, werden im nächsten Abschnitt geschildert.

\section{Ergebnisse}
Die Ergebnispräsentation zerfällt in eine allgemeine Darlegung der Ergebnisse des deskriptiven Intensivierers \textit{sehr} und der expressiven Intensivierer sowie in eine Begutachtung der Ergebnisse im Hinblick auf die zur Untersuchung stehenden Hypothesen H1 und H2 nach \citet[2]{Hauschild.1899} und \citet[274]{Baumgarten.1908}. Auf dieser Grundlage lassen sich die expressiven Intensivierer gemäß ähnlicher Frequenzen in Gruppen einteilen, deren Zusammensetzung eingehend analysiert wird. Wie zuvor erwähnt, beruhen die ermittelten Häufigkeiten auf den nach der jährlichen Korpusgröße bemessenen normalisierten Werten, sodass trotz eigentlicher Schwankungen diachrone Vergleichbarkeit gewährleistet ist.

Die statistischen Auswertungen und Grafiken wurden mithilfe der Statistik-Software \textsc{R} (Version 1.1.456, \citealt[]{R.2018}) und den Paketen \textit{ggplot2} (\citealt[]{Wickham.2016}) und \textit{dplyr} (\citealt[]{Wickham.2019}) durchgeführt und erstellt. Aufgrund der niedrigen Stichprobengröße wurde in dieser Untersuchung lediglich eine deskriptivstatistische Datenanalyse vorgenommen und von einer inferenzstatistischen Datenanalyse abgesehen.

\subsection{Deskriptiver Intensivierer}   \label{DeSehr}
\begin{figure}[b]
  \includegraphics[width=.8\textwidth]{sehr_diachron}
  \caption[Diachrone Entwicklung des deskriptiven Intensivierers \textit{sehr} im Zeitraum von 1950 bis 2017] {Diachrone Entwicklung des deskriptiven Intensivierers \textit{sehr} im Zeitraum von Januar 1950 bis Dezember 2017 in 5-Jahresintervallen (\mbox{x-Achse} = linearer Zeitverlauf, y-Achse = relative Häufigkeit).}
  \label{Dia-sehr}
\end{figure}

Wie infolge des hohen Lexikalisierungsgrades und der darauf zurückzuführenden stark ausgeprägten Verwendungsfreiheit \citep[vgl.][420]{Breindl.2007} zu erwarten war, ist das deskriptive \textit{sehr} der mit Abstand am häufigsten verwendete Intensivierer im Zeitschriftenkorpus. Es lässt sich über den Untersuchungszeitraum mit insgesamt 20\,278,12 Treffern pro Million Textwörter nachweisen, was verdeutlicht, wie unproblematisch sein Gebrauch auch in konzeptionell schriftlichen Textsorten ist. Dieses Ergebnis stimmt mit dem Befund von \citet[266]{Ito.2003} überein, die in ihrer Korpusuntersuchung für das englische \textit{very} dieselbe Beobachtung machten. Die diachronen Häufigkeitsverteilungen für \textit{sehr} sind in \figref{Dia-sehr} dargestellt.

Dem darin abgebildeten Verlauf nach zu urteilen, fällt die Verwendung von \textit{sehr} von Mitte der Fünfzigerjahre bis zu den Neunzigerjahren hin sichtbar ab. Seinen Tiefpunkt hat der deskriptive Intensivierer im Intervall um 1990 erreicht, wohingegen sein Gebrauch seit Mitte der Neunzigerjahre wieder verstärkt zunimmt. Es zeichnet sich diachron also ein Jo-Jo-Effekt mit abnehmendem Beginn, leichter Plateau-Phase und kontinuierlich zunehmendem Ende ab. Im letzten Intervall hat sich der Gebrauch von \textit{sehr} auf einem vergleichbar hohen Niveau wie im Zeitraum von 1955 bis 1960 eingependelt. Auffällig ist, dass \textit{sehr} trotz -- oder vielmehr aufgrund -- vollständiger Delexikalisierung, Konventionalisierung und fehlender Expressivität weiterhin stark im Gebrauch ist, was zweifellos seiner textsortenadäquaten skalaren Bedeutung geschuldet ist. Dies betrifft vermutlich nicht nur das zugrunde liegende Zeitschriftenkorpus, sondern mit hoher Wahrscheinlichkeit auch die Schriftsprache im Allgemeinen. Angesichts der hohen Frequenzen und des stabilen Vorkommens lassen sich bei \textit{sehr} keinerlei konkrete Anhaltspunkte für einen allmählichen Rückgang oder gar eine Verdrängung durch andere, im Besonderen expressive Intensivierer erkennen. Auch der Versuch, den Tiefpunkt zumindest partiell mit der Frequenz der expressiven Intensivierer zu korrelieren, misslang. Demnach findet sich keine korpuslinguistische Evidenz dafür, dass \textit{sehr} in seiner rückläufigen Phase, d.\,h. im Zeitraum von 1955 bis 1990, durch andere der hier untersuchten Ausdrücke ersetzt wurde. Insofern sind die möglichen Ursachen hierfür mehr auf anderen, eventuell auch außersprachlichen Ebenen zu suchen. Diese werden in Abschnitt \ref{DisKor} näher beleuchtet. Ebenso ist nicht auszuschließen, dass die Varianz statistisch nicht signifikant und daher ohnehin zu vernachlässigen ist.

\subsection{Expressive Intensivierer}
Die untersuchten expressiven Intensivierer \textit{arg}, \textit{derb}, \textit{furchtbar}, \textit{fürchterlich}, \textit{hammer}, \textit{krass}, \textit{mega}, \textit{ober}, \textit{sau}, \textit{scheiß}, \textit{schrecklich}, \textit{super}, \textit{unheimlich}, \textit{verdammt}, \textit{voll} und \textit{wahnsinnig} (n = 16) lassen sich im Zeitschriftenkorpus mit insgesamt 1\,719,36 Treffern pro Million Textwörter nachweisen (vgl. diesbezüglich auch die Übersicht in Tabelle \ref{tab:HäuVer} im Anhang). Somit ist der deskriptive Intensivierer \textit{sehr} mit 20\,287,12 Belegen mehr als zehn Mal so häufig im Korpus repräsentiert wie alle expressiven Intensivierer zusammen. Dies bekräftigt die im Vorfeld getroffene Vermutung, nach der \textit{sehr}, ähnlich wie bei der von \citet[]{Ito.2003} durchgeführten Korpusanalyse zu englischen Intensivierern, der meistverwendete Intensivierer in den beiden Zeitschriften ist. Offen bleibt jedoch, in welchen Textsorten die gesuchten Ausdrücke aufgetreten sind. Eine solche Einschätzung war aufgrund unzureichender
textexterner Metadaten und fehlender Textsortendokumentation nicht möglich. Ungeachtet dessen scheint der diachrone Gebrauch der expressiven Intensivierer im Korpus stetig angestiegen zu sein, wobei die Zunahme in den letzten fünf Dekaden besonders stark ausgeprägt ist, wie das in \figref{Dia-rest} illustrierte Säulendiagramm demonstriert.

\begin{figure}

  \includegraphics[width=.8\textwidth]{EI_diachron_neu.pdf}
  \caption[Diachrone Entwicklung der expressiven Intensivierer im Zeitraum von 1950 bis 2017] {Diachrone Entwicklung der expressiven Intensivierer (n = 16) im Zeitraum von Januar 1950 bis Dezember 2017 in 5-Jahresintervallen (\mbox{x-Achse} = linearer Zeitverlauf, y-Achse = relative Häufigkeit).}
  \label{Dia-rest}
\end{figure}

\noindent Der in \figref{Dia-rest} dargestellte Entwicklungsverlauf deutet auf eine positive Korrelation zwischen den Faktoren Zeit und Frequenz hin \citep[vgl.][105]{Perkuhn.2012}. So hat sich der im ersten 5-Jahresintervall gemessene IpMW-Wert von 15 bis zum letzten Intervall auf einen Wert von etwa 45 verdreifacht. Dies ist unter Berücksichtigung der eher konservativen Textsorten, aus denen sich eine Zeitschrift zusammensetzt, bemerkenswert.

Bei der Einzelbetrachtung lässt sich des Weiteren eine große Variabilität zwischen den einzelnen Ausdrücken ablesen, deren Vorkommen im Untersuchungskorpus verschieden stark ausgeprägt ist. Um dies zu veranschaulichen, sind die Entwicklungsverläufe der expressiven Intensivierer in \figref{SyDia} detailliert aufgeschlüsselt. Die Anordnung der Intensivierer erfolgt nach alphabetischer Reihenfolge. Zu beachten sind dabei die variierenden Skalen, die für die IpMW im Korpus stehen: Da die relative Frequenz der Ausdrücke stark variierte, ist die y-Achse jeweils unterschiedlich skaliert.

\begin{figure}
  \includegraphics[width=\textwidth]{EI_facets_neu3.pdf}
  \caption [Synchrone und diachrone Entwicklung der expressiven Intensivierer im Zeitraum von 1950 bis 2017] {Synchrone und diachrone Entwicklung der expressiven Intensivierer (n = 16) im Zeitraum von Januar 1950 bis Dezember 2017 in 3-Jahresintervallen (\mbox{x-Achse} = linearer Zeitverlauf, y-Achse = relative Häufigkeiten).}
  \label{SyDia}
\end{figure}

\noindent \figref{SyDia} weist auf unterschiedliche Entwicklungsmuster der expressiven Intensivierer hin. So gibt es Ausdrücke, die über die Zeit quantitativ eher zunehmen (z.\,B. \textit{wahnsinnig}), abnehmen (\textit{derb}) oder eine schwankende, nahezu wellenartige Entwicklung anzeigen (z.\,B. \textit{arg}). Um die diachronen Verläufe der Intensivierer im Detail zu inspizieren, werden anschließend Gruppen erstellt. Die Gruppeneinteilung wird dabei im Sinne der Hypothesen -- hier insbesondere H1 -- vorgenommen. Das bedeutet, dass bei der Einteilung sowohl dem Erstvorkommen (vgl. H1-Temp) wie auch der Korpusfrequenz (vgl. H1-Freq) eine tragende Rolle zukommt. Weil der Untersuchungszeitraum von 1950 bis 2017 etwa sieben Dekaden umfasst, wird er in drei Zeitabschnitte gegliedert, in denen die Ausdrücke quantitativ eingeordnet und miteinander verglichen werden:

\newpage
\begin{enumerate}
\item Zeitraum von 1950 bis 1969
\item Zeitraum von 1970 bis 1989
\item Zeitraum von 1990 bis 2017
\end{enumerate}

\noindent Somit umfasst jeder dieser Abschnitte 20 Jahre, wobei letzterem darüber hinaus auch die restlichen acht Jahre zugeschrieben werden. Wie \citet[97]{Perkuhn.2012} erklären, gibt es keine allgemeingültige Systematik hinsichtlich der Segmentierung von Korpora, was in Anbetracht der Variation von Korpusuntersuchungen wenig überraschend ist. Dementsprechend müssen korpuslinguistisch Forschende stets individuelle Segmentierungen unter Berücksichtigung ihrer Datengrundlage und Zielsetzung vornehmen. Um der verfolgten Fragestellung Rechnung zu tragen, ist meiner Ansicht nach eine Einteilung des Untersuchungskorpus in 20-Jahresschritte aus verschiedenen Gründen eine zielführende Vorgehensweise: Einerseits braucht es einige Zeit, bis sich typischerweise mündlichkeitsnahe Phänomene wie Intensitätspartikeln in der als eher konservativ geltenden Zeitungssprache etablieren. Aus diesem Grund sollte die Zeitspanne nicht zu knapp bemessen sein, wenn man valide Aussagen über die Entwicklung der Ausdrücke treffen möchte. Andererseits sollte die Zeitspanne jedoch auch nicht zu umfassend gewählt sein, da man sonst Gefahr läuft, potenziell bedeutsame Entwicklungsschritte zu übersehen oder zu verallgemeinern. Schaut man sich die Entwicklungsverläufe der expressiven Intensivierer in den Abbildungen \ref{Dia-rest} und \ref{SyDia} an, so scheinen Intervalle von 20 Jahren eine gute Grundlage zu sein, anhand derer der Fortgang der Ausdrücke in Bezug auf die Hypothesen abgeglichen werden kann. So bildet das Ende dieses Zeitraums oftmals die Trennlinie für offensichtliche Veränderungen, wie bspw. an dem um 1970 herum ausgeprägten Tiefpunkt von \textit{arg} und \textit{derb} zu erkennen ist.

Neben der zeitlichen Dimension, die Rückschlüsse auf das Erstvorkommen der Ausdrücke ermöglicht und überdies erste Hinweise auf das Zutreffen der Hypothese H2 liefert, wird die Häufigkeit der Intensivierer pro Segment begutachtet. In diesem Kontext stellt sich die Frage, welche expressiven Intensivierer in den Zeitabschnitten hoch- und welche niedrigfrequent sind. Aus Gründen der Vollständigkeit wird \textit{sehr} als Vertreter der deskriptiven Intensiviererklasse ebenso einer solchen Begutachtung unterzogen, wenn auch die Expressivität hier nicht als Erklärung für die zu beobachtenden Verläufe in Erwägung gezogen wird. Aufgrund dessen, dass es sich um einen einzelnen Ausdruck handelt, wird für \textit{sehr} zunächst jede Dekade separat aufgeführt, woraus sich schließlich die kumulierten Werte für die oben erläuterten Zeitabschnitte (d.\,h. 1950 bis 1969, 1970 bis 1989 und 1990 bis 2017) ergeben.

Ehe die Ergebnisse der Auswertung vorgestellt werden, sei darauf hingewiesen, dass die Grenzen keineswegs als trennscharf zu beurteilen sind. Dennoch lassen die Tendenzen bedeutsame Mutmaßungen über den Gebrauch der Ausdrücke (stellvertretend für konventionell schriftliche Kontexte) zu.

\subsection[Segmentale Entwicklung des deskriptiven Intensivierers]{Segmentale Entwicklung des deskriptiven Intensivierers}
Wie in Abschnitt \ref{DeSehr} geschildert, lässt sich die diachrone Entwicklung von \textit{sehr} als eine Art Welle mit abnehmendem Beginn, verhältnismäßig niedrigerer Plateau-Phase und ansteigendem Ende beschreiben.

Aus diesem Grund ist es interessant, sich die kumulierte Häufigkeitsverteilung des deskriptiven Intensivierers über den Untersuchungszeitraum von 1950 bis 2017 genauer anzuschauen. Die Frequenzen von \textit{sehr} je Dekade sind in chronologischer Reihenfolge in Tabelle \ref{Fre0} zusammengetragen.

%%%%%%%%%%%

\begin{table}
\fittable{
\begin{tabular}{l@{~}rrc@{}r@{~}}
\lsptoprule
{Jahrzehnt} & Absolute  Frequenz  & Relative Frequenz  & & Kumulierte Frequenz \\
\midrule
1950--1959 & 12\,638 & 3\,522,92 & \rdelim\}{2}{.4cm} & {\multirow{2}{*} {6\,714,94}}\\
1960--1969 & 20\,949 & 3\,192,02 &  & \\
\tablevspace
1970--1979 & 21\,192 & 2\,692,41 & \rdelim\}{2}{.4cm} & {\multirow{2}{*} {5\,190,35}}\\
1980--1989 & 24\,662 & 2\,497,94 &  & \\
\tablevspace
1990--1999 & 21\,770 & 2\,605,88 &  \rdelim\}{3}{.4cm} & {\multirow{3}{*} {8\,372,83}} \\
2000--2009 & 34\,632 & 2\,916,41 &  & \\
2010--2017 & 29\,318 & 2\,850,54 &  & \\
\midrule
In Summe & 165\,161 & 2\,0278,12 &  &\\
\lspbottomrule
\end{tabular}
}
\caption[Kumulierte Häufigkeitsverteilung des deskriptiven Intensivierers im Zeitraum von 1950 bis 2017]{Kumulierte Häufigkeitsverteilung des deskriptiven Intensivierers \textit{sehr} im Zeitraum von 1950 bis 2017 (in chronologischer Reihenfolge).}
\label{Fre0}
\end{table}

\noindent Wie das in \figref{Dia-sehr} dargestellte Histogramm vermuten lässt, geben auch die Korpusfrequenzen von \textit{sehr} eine zunächst rückläufige Entwicklung an. So gehen die in den ersten beiden Dekaden ermittelten Frequenzen von über 3\,000 Treffern sukzessive derart zurück, dass der Intensivierer im Zeitraum von 1980 bis 1989 mehr als 1\,000 Korpusbelege weniger als in der ersten Dekade aufweist. Nach dem Tiefpunkt nimmt \textit{sehr} in seiner Frequenz schließlich wieder zu, sodass das Intervall von 2000 bis 2009 vergleichbar mit dem zweiten Intervall von 1960 bis 1969 ist. Da der letzte Zeitabschnitt aufgrund des beschränkten Untersuchungszeitraums lediglich bis Dezember 2017 reicht und somit noch 24 Monate bis zur Vollendung der siebten Dekade ausstehen, lassen sich die fehlenden Werte an dieser Stelle nur extrapolieren. Ausgehend von der Vermutung, dass die Verwendung von \textit{sehr} in den Jahren 2018 und 2019 in etwa gleich stark ausgeprägt ist, ist zu erwarten, dass das kumulierte Vorkommen in diesem Jahrzehnt dem aus dem ersten weitgehend in Nichts nachsteht. Dies lässt darauf schließen, dass der deskriptive Intensivierer seine quantitative Einbuße überstanden und sich der Gebrauch im Korpus nach temporärem Rückgang bis zur vierten Dekade schließlich zum Positiven gewendet hat. Gegenwärtig liegt der Gebrauch von \textit{sehr} im Zeitschriftenkorpus in etwa auf einem vergleichbar hohen Frequenzniveau wie zu Beginn des Untersuchungszeitraums.

\subsection{Segmentale Entwicklung der expressiven Intensivierer}
Anders als bei \textit{sehr} erfolgt die Begutachtung der expressiven Intensivierer direkt separiert nach den oben umschriebenen Zeitspannen. Dies gewährt zum einen einen besseren Überblick über die Häufigkeitsverteilungen. Zum anderen ermöglicht es auch eine zielgerichtetere Gegenüberstellung der Ausdrücke. Auf Basis der bezüglich H1 relevanten Faktoren Erstvorkommen (vgl. H1-Temp) und Frequenz (vgl. H1-Freq) werden die expressiven Intensivierer sukzessive hierarchisch klassifiziert. Im Zuge dessen wird im ersten Analyseschritt für jeden Zeitabschnitt überprüft, welche der Ausdrücke im Korpus alt und welche neu sind. Exemplarisch wird am zweiten Intervall daher der Frage nachgegangen, welche Ausdrücke es schon vor 1969 gab und welche erst anschließend erstmals in intensivierender Funktion in Erscheinung treten. Im zweiten Analyseschritt wird die Frequenz beleuchtet: So lassen sich in jedem Intervall zwei Frequenzgruppen ausmachen, deren Zusammensetzung eingehend in Augenschein genommen wird.

Unterstützt wird diese erste Einteilung durch eine mit R vorgenommene hierarchische agglomerative Clusteranalyse, die mit dem Ziel durchgeführt wurde, durch die daraus abgeleiteten Dendrogramme Ähnlichkeiten zwischen den Intensivierern aufzudecken. Dabei handelt es sich um ein exploratives Bottom-up-Verfahren, mit dem es auf Grundlage der Ergebnisse möglich ist, Hypothesen für weitere Untersuchungen aufzustellen \citep[vgl.][294]{Gries.2008}. Auf diese Weise lässt sich mithilfe einer grafischen Visualisierung einschätzen, ob die erste, intuitiv unternommene Gruppeneinteilung treffend ist oder revidiert bzw. \mbox{justiert} werden muss. Bei einer solchen Analyse wird eine Anzahl von Objekten (hier: Intensivierer) gemäß bestimmter Merkmale (hier: Frequenzen) in verschiedene Cluster eingeteilt. Die Eingruppierung der Ausdrücke wird durch das relative Vorkommen der Ausdrücke im Korpus, d.\,h. die IpMW, generiert. Der Ausgangspunkt liegt bei 16 Einzelgruppen, alle expressiven Intensivierer also, die mittels einer Distanzmatrix zu homogenen Gruppen zusammengefasst werden. Die Eigenschaften der Mitglieder eines Clusters ähneln sich, während sie sich von den Mitgliedern anderer Gruppen hingegen unterscheiden. Demzufolge sind die Intensivierer innerhalb eines Clusters hinsichtlich der Korpusfrequenz homogen und heterogen zu den Intensivierern der anderen Cluster. Dabei gilt: Je niedriger die Distanz zwischen den einzelnen Gruppen ist, desto ähnlicher sind sich die Ausdrücke. Oder andersherum ausgedrückt: Je größer die Distanz zwischen den einzelnen Gruppen ist, desto unähnlicher sind sie sich im Hinblick auf ihre Korpusfrequenz (vgl. \mbox{\citealt[296]{Gries.2008}}; \mbox{\citealt[440]{Backhaus.2018}}; \mbox{\citealt[437ff.]{Sauer.2019}}).

Die für die jeweiligen Zeitabschnitte abgeleiteten Dendrogramme befinden sich unter den Ausführungen zu den Frequenzen eines jeden Intervalls. In diesen sind die expressiven Intensivierer in Form einer grafischen Baumstruktur auf Grundlage der Euklidischen Distanz\footnote{Die Euklidische Distanz ist das Maß, das die Ähnlichkeit bzw. Unähnlichkeit der Intensivierer auf Grundlage ihrer Entfernung voneinander numerisch ausdrückt. Nach \citet[450]{Backhaus.2018} werden dabei \glqq{}für jedes Objektpaar die Differenzwerte jeder Eigenschaft quadriert und anschließend addiert\grqq{}, woraufhin aus der Summe die Quadratwurzel gezogen wird.} in verschiedene, homogene Cluster eingruppiert. Die x-Achse symbolisiert die Distanz, auf deren Basis sich Rückschlüsse in Bezug auf die Gemeinsamkeiten bzw. Unterschiede der Ausdrücken ziehen lassen, die y-Achse steht für den Grad der Ähnlichkeit der Ausdrücke und repräsentiert das Fusionskriterium. Die Verschmelzung der Intensivierer erfolgte durch das Ward-Verfahren (\citealt[]{Ward.1963}),\footnote{Anders als andere Methoden gruppiert das Ward-Verfahren als varianzanalytische Amalgamierungsregel Objekte, die keinen großen Einfluss auf ein vorgegebenes Heterogenitätsmaß bzw. Varianzkriterium ausüben. Dadurch wird die Streuung innerhalb einer Gruppe so wenig wie möglich erhöht und die Fehlerquadratsumme entsprechend niedrig gehalten {(vgl. \citealt[305]{Gries.2008}; \citealt[465f.]{Backhaus.2018}}).} das nach möglichst homogenen Clustern strebt \citep[vgl.][305]{Gries.2008}. Nachdem das genaue Vorgehen transparent gemacht wurde, folgt im nächsten Part die Betrachtung des ersten Zeitabschnitts, d.\,h. des Zeitraums von 1950 bis 1969.

\subsubsection{Zeitabschnitt von 1950 bis 1969}
Im ersten Zeitabschnitt finden sich im Korpus 13 expressive Intensivierer, die mit insgesamt 277,76 Treffern pro Million Textwörter auftreten. Begonnen wird jeweils mit denjenigen Ausdrücken, die eine niedrige Frequenz aufweisen. Dies betrifft im ersten Intervall \textit{scheiß}, \textit{voll}, \textit{ober}, \textit{fürchterlich}, \textit{wahnsinnig}, \textit{sau} und \textit{derb} (n = 7). Die Gruppe der niedrigfrequenten Intensivierer bewegt sich quantitativ in einem Rahmen von 0 bis ca. 12 Belegen (${\scriptstyle \sum{}}$ = 41,53)\footnote{Das Symbol Sigma (${\scriptstyle \sum{}}$) meint hier und im Folgenden die Summe der relativen Frequenzen.}. Die zweite Gruppe beinhaltet dagegen diejenigen Ausdrücke, deren relative Frequenz im Korpus bei mehr als 25 Treffern liegt (${\scriptstyle \sum{}}$ = 236,23): \textit{unheimlich}, \textit{furchtbar}, \textit{arg}, \textit{verdammt}, \textit{schrecklich} und \textit{super} (n = 6). Zur Illustration sind die beiden Gruppen mitsamt der jeweiligen kumulierten Frequenz der Mitglieder in Tabelle \ref{Fre1} abgebildet. Die Sortierung der Intensivierer je Gruppe erfolgt in zunehmender Frequenz.

%%%%%%%%%%%%%%%%%

\begin{table}
\begin{tabularx}{\textwidth}{l@{}>{\raggedleft}p{2cm}@{}>{\raggedleft}p{2cm}@{}>{\raggedleft}p{2.5cm}@{~}r@{}Y}
\lsptoprule
{Intensivierer} & {Absolute} Frequenz & {Relative} Frequenz & {Prozentualer} Anteil & & {Klassifikation} \\
\midrule
\textit{scheiß} & {7} & {1,18} & {0,42\,\%} & \rdelim\}{7}{.3cm} & \\
\textit{voll} & {9} &{1,68} & {0,60\,\%} & &\\
\textit{ober} & {11} &{3,01} & {1,08\,\%} & & \multirow{3}{*}{\parbox{2.5cm}{{niedrigfrequent \\ (14,95\,\%)}}} \\
\textit{fürchterlich} & {37} &{6,27} & {2,26\,\%} & &  \\
\textit{wahnsinnig} & {29} & {6,90} & {2,48\,\%} & & \\
\textit{sau} & {49} &{10,78} & {3,88\,\%} & & \\
\textit{derb} & {55} &{11,71} & {4,22\,\%} & &\\
\midrule
\textit{unheimlich} & {131} &{25,69} & {9,25\,\%} & \rdelim\}{6}{.3cm} & \\
\textit{furchtbar} & {173} & {35,73} & {12,86\,\%} & &  \multirow{4}{*}{\parbox{2.5cm}{{hochfrequent \\ (85,05\,\%)}}}\\
\textit{arg} & {191} &{38,21} & {13,76\,\%} & & \\
\textit{verdammt} & {195} &{39,94} & {14,38\,\%} & & \\
\textit{schrecklich} & {194} &{40,65} & {14,63\,\%} & &\\
\textit{super} & {263} &{56,01} & {20,16\,\%} & &\\
\midrule
In Summe & {1\,344} & {277,76} & {100\,\%} & &\\
\lspbottomrule
\end{tabularx}
\caption[Kumulierte Häufigkeitsverteilungen der expressiven Intensivierer im Zeitraum von 1950 bis 1969]{Kumulierte Häufigkeitsverteilungen der expressiven Intensivierer im Zeitraum von 1950 bis 1969 (angeordnet nach zunehmender Frequenz).}
\label{Fre1}
\end{table}

\noindent Wie sich aus Tabelle \ref{Fre1} ablesen lässt, verfügt \textit{unheimlich} über die niedrigste Trefferquote von Gruppe 2. Dennoch ist dieser Wert mehr als doppelt so hoch wie die Trefferzahl von \textit{derb}, das in Gruppe 1 mit 11,71 Treffern pro Million Textwörter die höchste Frequenz aufweist. Infolge dieser recht ausgeprägten Differenz scheint eine Grenze an dieser Stelle, d.\,h. zwischen \textit{derb} und \textit{unheimlich}, gerechtfertigt. Interessanterweise liegen die Intensivierer \textit{furchtbar}, \textit{arg}, \textit{verdammt} und \textit{schrecklich}, die qua Semantik einen negativen Zustand denotieren, zahlenmäßig eng beieinander. Demgegenüber ist \textit{super} in diesem Intervall mit 56 Korpusbelegen sogar doppelt so häufig in den beiden Zeitschriften zu finden wie \textit{unheimlich}. Auffällig ist weiterhin auch, dass die Intensivierer \textit{hammer}, \textit{mega} und \textit{krass} im ersten Zeitabschnitt überhaupt nicht vorkommen. Bei ihnen handelt es sich also um Ausdrücke, die sich erst innerhalb des Untersuchungszeitraums als Intensivierer im Korpus neu herausgebildet haben.

In \figref{Dendro1} ist das für den Zeitraum von 1950 bis 1969 einschlägige Dendrogramm dargestellt. Anhand diesem lässt sich überprüfen, ob die in der obigen Tabelle skizzierte Gruppeneinteilung zutreffend ist.

\begin{figure}

\includegraphics[width=1\textwidth]{1950-1969}
\caption[Cluster-Dendrogramm der expressiven Intensivierer für den Zeitraum von 1950 bis 1969]{Cluster-Dendrogramm der expressiven Intensivierer für den Zeitraum von 1950 bis 1969 (Maß: Euklidische Distanz; Amalgamierungsmethode: Ward-Verfahren (\citealt[]{Ward.1963})).}
\label{Dendro1}
\end{figure}

\noindent Das Dendrogramm aus \figref{Dendro1} zeigt zwei große Cluster mit mehreren Subclustern. Das linke Cluster beinhaltet dabei sowohl diejenigen Intensivierer, die sich in diesem Zeitabschnitt im Korpus überhaupt nicht nachweisen lassen (d.\,h. \textit{mega}, \textit{hammer} und \textit{krass}), als auch solche, die mit einer niedrigen Frequenz auftreten: \textit{derb}, \textit{sau}, \textit{fürchterlich}, \textit{wahnsinnig}, \textit{voll}, \textit{ober} und \textit{scheiß}. Intensivierer, die über homogene Merkmale verfügen, also  ähnlich frequent sind, werden zudem in weiteren Subclustern zusammengefasst, was u.\,a. am Beispiel von \textit{derb} und \textit{sau} zu sehen ist. Auf der anderen Seite finden sich im rechten Cluster diejenigen Intensivierer, die im Korpus deutlich frequenter sind, nämlich \textit{super}, \textit{unheimlich}, \textit{furchtbar}, \textit{arg}, \textit{schrecklich} und \textit{verdammt}. Wie bereits im Hinblick auf Tabelle \ref{Fre1} angesprochen, beinhaltet diese Gruppe mit Ausnahme von \textit{super} ausschließlich Ausdrücke mit klar negativer Quellsemantik, die die durch das Adjektiv ausgedrückte Eigenschaft in Form von Furcht, Angst und Schrecken verstärken (vgl. \citealt[165ff.]{Biedermann.1969}; \citealt[]{Hentschel.1998b}). Aufgrund vergleichbarer Frequenzen liegen diese Ausdrücke auch im Dendrogramm nahe beieinander. Wie das Fusionskriterium nahelegt, unterscheidet sich der Intensivierer \textit{super} stark von den restlichen Gruppenmitgliedern, was dessen ungleich höherer Korpusfrequenz geschuldet ist. Das Dendrogramm des ersten Zeitabschnitts scheint also die zuvor angenommene Gruppeneinteilung aus Tabelle \ref{Fre1} vollständig zu stützen, sodass keine Nachjustierungen erforderlich sind.

\subsubsection{Zeitabschnitt von 1970 bis 1989}
Im zweiten Intervall sind alle 16 expressiven Intensivierer mit einer kumulierten Gesamtfrequenz von 485,56 Treffern im Zeitschriftenkorpus vertreten. Davon verfügen die Ausdrücke \textit{krass}, \textit{hammer}, \textit{mega}, \textit{ober}, \textit{voll}, \textit{derb}, \textit{sau}, \textit{scheiß} und \textit{fürchterlich} (n = 9) über ein niedriges Vorkommen. Diese weisen im Zeitraum von 1970 bis 1989 maximal 14 Korpusbelege auf (${\scriptstyle \sum{}}$ = 53,41). Auf der anderen Seite finden sich hingegen die Intensivierer \textit{wahnsinnig}, \textit{furchtbar}, \textit{verdammt}, \textit{schrecklich}, \textit{unheimlich}, \textit{arg} und \textit{super} (n = 7), die mit mindestens 26 Treffern merklich frequenter sind (${\scriptstyle \sum{}}$ = 432,15). Wie zuvor sind die genauen Korpusfrequenzen im Sinne der Übersichtlichkeit tabellarisch dargestellt (vgl. Tabelle \ref{Fre2}). Auch hier erfolgt die Anordnung jeweils nach zunehmender Frequenz.

\begin{table}
\begin{tabularx}{\textwidth}{l@{}>{\raggedleft}p{2cm}@{}>{\raggedleft}p{2cm}@{}>{\raggedleft}p{2.5cm}@{~}r@{}Y}
\lsptoprule
{Intensivierer} & {Absolute}  Frequenz& {Relative} Frequenz& {Prozentualer} Anteil& & {Klassifikation} \\
\midrule
\textit{krass} & {1} & {0,10} & 	{0,02\,\%}		& \rdelim\}{9}{.3cm} & \\
\textit{hammer} & {1} & {0,12} &	{0,02\,\%}		& 			&\\
\textit{mega} & {2} &{0,24}& 		{0,05\,\%}	& &\\
\textit{ober} & {27} & {2,83} &		{0,58\,\%}	& &  \multirow{3}{*}{\parbox{2.5cm}{{niedrigfrequent \\ (11,02\,\%)}}} \\
\textit{voll} & {29} & {3,15} &		{0,65\,\%}	& & \\
\textit{derb} & {52} &{6,00} 	&	{1,24\,\%}	& &  \\
\textit{sau} & {121} & {13,41} &		{2,77\,\%}	& &\\
\textit{scheiß} & {123} & {13,71} &	{2,83\,\%}	& &\\
\textit{fürchterlich} & {122} & {13,85} &	{2,86\,\%} & &\\
\midrule
\textit{wahnsinnig} & {239} & {26,15} & {5,19\,\%} & \rdelim\}{7}{.3cm} & \\
\textit{furchtbar} & {403} & {44,99} & {9,28\,\%} & &\\
\textit{verdammt} & {460} & {51,26} &  {10,58\,\%} & &  \multirow{3}{*}{\parbox{2.5cm}{{hochfrequent \\ (88,98\,\%)}}} \\
\textit{schrecklich} & {508} & {56,55}& {11,67\,\%} & & \\
\textit{unheimlich} & {579} & {64,23} & {13,26\,\%} & & \\
\textit{arg} & {602} & {66,85}& {13,80\,\%} & &\\
\textit{super}  & {1\,075} &  {122,12}& {25,20\,\%} & &\\

\midrule
In Summe & {4\,344} & {485,56} & {100\,\%} &  &\\
\lspbottomrule
\end{tabularx}
\caption[Kumulierte Häufigkeitsverteilungen der expressiven Intensivierer im Zeitraum von 1970 bis 1989]{Kumulierte Häufigkeitsverteilungen der expressiven Intensivierer im Zeitraum von 1970 bis 1989 (angeordnet nach zunehmender Frequenz).}
\label{Fre2}
\end{table}

%%%%%%%%%%%%%%%%%%

Da \textit{wahnsinnig} im Korpus beinahe doppelt so frequent ist wie \textit{fürchterlich}, erscheint mir eine Grenze an dieser Stelle sinnvoll. Die im Korpus neu entwickelten Intensivierer \textit{krass}, \textit{mega} und \textit{hammer} weisen in diesem Zeitabschnitt eine nur geringe Frequenz auf. Daneben hat sich innerhalb von Gruppe 1 die häufigkeitsbedingte Anordnung geändert, sodass \textit{fürchterlich} hier nun am häufigsten ist. Der im vorigen Zeitabschnitt niedrigfrequente Intensivierer \textit{wahnsinnig} ist in diesem Intervall in die Gruppe der hochfrequenten Intensivierer übergetreten, sodass diese auf insgesamt sieben Mitglieder angewachsen ist. Darin lässt sich \textit{super} als einziger Intensivierer mit mehr als 100 Belegen nachweisen, sodass er seine Spitzenposition stark ausgebaut hat. Bei den restlichen Ausdrücken dieser Gruppe handelt es sich mit \textit{wahnsinnig}, \textit{furchtbar}, \textit{verdammt}, \textit{schrecklich}, \textit{unheimlich} und \textit{arg} noch immer ausschließlich um solche, die sich durch eine transparente negative Semantik auszeichnen. Auch diese Intensivierer sind im Vergleich zum vorangegangenen Intervall in einer anderen Reihenfolge angeordnet.


\begin{figure}

\includegraphics[width=1\textwidth]{1970-1989}
\caption[Cluster-Dendrogramm der expressiven Intensivierer für den Zeitraum von 1970 bis 1989]{Cluster-Dendrogramm der expressiven Intensivierer für den Zeitraum von 1970 bis 1989 (Maß: Euklidische Distanz; Amalgamierungsmethode: Ward-Verfahren (\citealt[]{Ward.1963})).}
\label{Dendro2}
\end{figure}

In \figref{Dendro2} ist das Dendrogramm für den Zeitabschnitt 1970 bis 1989 dargestellt. Dieses zerfällt erneut in zwei Cluster, die sich wiederum aus mehreren Subclustern zusammensetzen. Auch hier befinden sich im linken Cluster Ausdrücke mit niedriger Korpusfrequenz: \textit{sau}, \textit{fürchterlich}, \textit{scheiß}, \textit{voll}, \textit{derb}, \textit{ober}, \textit{mega}, \textit{hammer} und \textit{krass}. Die neuen Intensivierer \textit{mega}, \textit{hammer} und \textit{krass} sind aufgrund ähnlichen Vorkommens im Korpus in ein Subcluster zusammengefasst. Im Gegensatz dazu beinhaltet das zweite Cluster die frequenteren Ausdrücke und wird dementsprechend durch \textit{super}, \textit{arg}, \textit{schrecklich}, \textit{unheimlich}, \textit{verdammt}, \textit{furchtbar} und \textit{wahnsinnig} gebildet. Im Dendrogramm ist weiterhin der beschriebene Wechsel des vormals niedrigfrequenten Intensivierers \textit{wahnsinnig} in die hochfrequente Gruppe zu sehen, genauso wie sich darin die stärkere Frequenzzunahme von \textit{super} herauskristallisiert. Gleichwohl die \glqq{}schrecklichen\grqq{} Intensivierer  im Korpus vergleichsweise frequenter sind, hat sich das Vorkommen von \textit{super} in diesem Intervall sogar mehr als verdoppelt. Auch dieses Dendrogramm deutet auf ein Zutreffen der für den Zeitabschnitt von 1970 bis 1989 vorgenommenen Gruppeneinteilung hin, ohne dass Anpassungen vonnöten sind.

\subsubsection{Zeitabschnitt von 1990 bis 2017} \label{abs:Freqgruppe}
Wie schon im vorherigen Intervall sind auch im letzten alle 16 expressiven Intensivierer im Zeitschriftenkorpus vertreten. Diese lassen sich mit insgesamt 956,04 Treffern pro Million Textwörter nachweisen. Die Gruppe der niedrigfrequenten Ausdrücke setzt sich wiederholt aus \textit{derb}, \textit{krass}, \textit{hammer}, \textit{mega}, \textit{ober}, \textit{fürchterlich}, \textit{scheiß}, \textit{voll} und \textit{sau} (n = 9) zusammen. Diese Intensivierer kommen im Korpus mit maximal 37 Belegen vor (${\scriptstyle \sum{}}$ = 163,30). Erheblich häufiger sind dagegen die Intensivierer \textit{unheimlich}, \textit{arg}, \textit{schrecklich}, \textit{furchtbar}, \textit{wahnsinnig}, \textit{verdammt} und \textit{super} (n = 7) zu finden, die auch in diesem Zeitabschnitt die Gruppe der hochfrequenten Intensivierer bilden (${\scriptstyle \sum{}}$ = 792,74) und im Korpus mit mindestens 72 Treffern nachweisbar sind. Die beiden Frequenzgruppen sind inklusive der kumulierten Häufigkeiten ihrer Mitglieder in Tabelle \ref{Fre3.1} detailliert aufgeschlüsselt. Die Sortierung der Gruppenmitglieder erfolgt in zunehmender Frequenz.

\begin{table}
\begin{tabularx}{\textwidth}{l@{}>{\raggedleft}p{2cm}@{}>{\raggedleft}p{2cm}@{}>{\raggedleft}p{2.5cm}@{~}r@{}Y}
\lsptoprule
{Intensivierer} & {Absolute}  Frequenz& {Relative} Frequenz& {Prozentualer} Anteil& & {Klassifikation} \\
\midrule
\textit{derb}  & {42} & {4,35}	&	{0,46\,\%}	& \rdelim\}{9}{.3cm} &  \\
\textit{krass} & {72} &	{7,01}	&{0,73\,\%}	& 			&\\
\textit{hammer} & {77} &	{7,58}&		{0,79\,\%}& &\\
\textit{mega} & {113} &	{11,19}	&	{1,17\,\%}& &  \multirow{3}{*}{\parbox{2.5cm}{{niedrigfrequent \\ (17,09\,\%)}}}  \\
\textit{ober} & {147} &	{14,69}	&	{1,54\,\%}& &  \\
\textit{fürchterlich} & {238} & {23,36}&	{2,44\,\%}		& &  \\
\textit{scheiß} & {273} & 	{27,08}&	{2,83\,\%}	& &\\
\textit{voll} & {319} & {31,82}	&	{3,33\,\%}& &\\
\textit{sau} & {353} & {36,22}&	{3,79\,\%} & & \\
\midrule
\textit{unheimlich} & {717} & {71,97}& {7,53\,\%} & \rdelim\}{7}{.3cm} & \\
\textit{arg} & {801} & {80,44}& {8,41\,\%} & &\\
\textit{schrecklich} & {898} & {89,06}& {9,32\,\%} & & \multirow{3}{*}{\parbox{2.5cm}{{hochfrequent \\ (82,91\,\%)}}}\\
\textit{furchtbar} & {981} & {96,39}& {10,08\,\%} & & \\
\textit{wahnsinnig} &  {1\,083} & {104,53}& {10,93\,\%}& &     \\
\textit{verdammt} & {1\,365} &  {137,22}&{14,35\,\%} & &\\
\textit{super} & {2\,138} & {213,13}&{22,29\,\%}  & & \\
\bottomrule   %82,91\\
In Summe & {9\,617} & {956,04} & {100\,\%} & &\\
\lspbottomrule
\end{tabularx}

\caption[Kumulierte Häufigkeitsverteilungen der expressiven Intensivierer im Zeitraum von 1990 bis 2017]{Kumulierte Häufigkeitsverteilungen der expressiven Intensivierer im Zeitraum von 1990 bis 2017 (angeordnet nach zunehmender Frequenz).}
\label{Fre3.1}
\end{table}

%%%%%%%%%%%%%%%%%%

\noindent In diesem Zeitraum weist \textit{unheimlich} im Vergleich zu \textit{sau} nahezu doppelt so viele Treffer auf, was für eine Grenze zwischen diesen beiden Ausdrücken spricht. Ebenso haben in diesem Intervall die rezent etablierten Intensivierer \textit{krass}, \textit{mega} und \textit{hammer} an Frequenz dazugewonnen und den älteren Ausdruck \textit{derb} hinter sich gelassen. Im Hinblick auf die Gruppenzusammensetzung zeigt sich verhältnismäßig wenig Veränderung. So gliedert sich die Gruppe der hochfrequenten Intensivierer noch immer vorwiegend in die negativ konnotierten Ausdrücke mit transparenter Quellsemantik. Lediglich \textit{super} stellt in diesem Punkt eine Ausnahme dar, stammt der Ausdruck doch aus dem Lateinischen und gibt qua Semantik an, dass eine kontextuell festgelegte Norm überschritten wird. Die kumulierten Frequenzen demonstrieren darüber hinaus, dass in diesem Zeitabschnitt ebenfalls die beiden Intensivierer \textit{wahnsinnig} und \textit{verdammt} den Wert von 100 Belegen überschritten haben, wohingegen \textit{super} seinen Vorsprung sogar auf über 200 Treffer pro Million Textwörter verdoppelt hat. Somit ist der Ausdruck auch im dritten Intervall der im Korpus am häufigsten vorkommende Intensivierer.

\figref{Dendro3} illustriert das Dendrogramm für die Zeit von 1990 bis 2017.

\begin{figure}

\includegraphics[width=1\textwidth]{1990-2017}
\caption[Cluster-Dendrogramm der expressiven Intensivierer für den Zeitraum von 1990 bis 2017]{Cluster-Dendrogramm der expressiven Intensivierer für den Zeitraum von 1990 bis 2017 (Maß: Euklidische Distanz; Amalgamierungsmethode: Ward-Verfahren (\citealt[]{Ward.1963})).}\label{Dendro3}
\end{figure}

\noindent Wie die vorigen setzt sich das letzte Dendrogramm aus zwei Clustern mit mehreren Subclustern zusammen. Das Cluster, das die niedrigfrequenten Intensivierer beinhaltet, zerfällt erwartungsgemäß in \textit{derb}, \textit{hammer}, \textit{krass}, \textit{mega}, \textit{ober}, \textit{fürchterlich}, \textit{scheiß}, \textit{sau} und \textit{voll}. In Bezug auf seine Häufigkeit unterscheidet sich \textit{derb} in diesem Zeitabschnitt erstmals stark von den restlichen niedrigfrequenten Ausdrücken. Außerdem haben auch die neuen Intensivierer sichtbare Entwicklungen durchlaufen: So bilden \textit{hammer} und \textit{krass} aufgrund ähnlicher Frequenzen ein Subcluster, während \textit{mega} zusammen mit \textit{ober} ein weiteres bildet. Demnach kommt es im dritten Intervall also erstmals zu einer Trennung der rezent in Erscheinung getretenen Ausdrücke, deren Ursache in den unterschiedlichen Frequenzverläufen liegt. Noch ausgeprägter sind die Unterschiede im zweiten Cluster, das die hochfrequenten Intensivierer \textit{super}, \textit{verdammt}, \textit{wahnsinnig}, \textit{furchtbar}, \textit{schrecklich}, \textit{arg} und \textit{unheimlich} umfasst. Obwohl sich bezüglich der Gruppenzusammensetzung im Vergleich zum vorherigen Zeitabschnitt per se keine Veränderungen ergeben haben, sind die Unterschiede zwischen den Subclustern hingegen erheblich markanter. Diese lassen sich auf die Zunahme der Korpusfrequenzen zurückführen, die eine klarere Ausdifferenzierung der Subgruppen nach sich zieht. Die Zusammensetzung der Subcluster unterscheidet sich dabei in dem Sinne von der des vorangegangenen Intervalls, dass in diesem Zeitraum neue frequenzbedingte Intensiviererpaare auszumachen sind, bspw. die beiden hochfrequenten Intensivierer \textit{arg} und \textit{unheimlich}. Noch auffälliger ist die Entwicklung, die \textit{super} vom zweiten zum dritten Zeitabschnitt durchlaufen hat: Wies es im zweiten Intervall noch ca. 122 Korpusbelege auf, so ist sein Gebrauch im dritten auf über 200 Treffer pro Million Textwörter angestiegen. Im Dendrogramm resultiert diese Entwicklung in einem ausgeprägteren Unterschied zwischen \textit{super} und der Gruppe der \glqq{}schrecklichen\grqq{} Intensivierer, der sich vom ersten bis zum dritten Zeitraum sukzessive gesteigert hat. Kurz gefasst deutet auch dieses Dendrogramm auf ein Zutreffen der für das Intervall von 1990 bis 2017 angenommenen Gruppeneinteilung hin, ohne dass Anpassungen vorzunehmen sind.

\subsubsection{Direkter Vergleich der Zeitabschnitte}
Nachdem die festgelegten Intervalle der Reihe nach analysiert wurden, folgt in diesem Abschnitt ein direkter Vergleich ebendieser, der sich nach den für die Hypothese H1 zentralen Faktoren Erstvorkommen (vgl. H1-Temp) und Frequenz (vgl. H1-Freq) richtet. Um einen besseren Überblick über die Gegenüberstellung zu geben, werden die kumulierten Häufigkeiten der Intensivierer pro Zeitraum inklusive ihrer prozentualen Entwicklungstendenz in Tabelle \ref{Fre4} komprimiert dargelegt. Diese  Werte sind in zunehmender Gesamtfrequenz gelistet. In der Tabelle sind jeweils diejenigen Spalten farblich hervorgehoben, bei denen sich eine diachrone Frequenzabnahme abzeichnet. Wie zu Beginn des Kapitels beschrieben, ist im Sinne der in Abschnitt \ref{Auslöser} umrissenen Expressivitätsmaximen nach \citet[]{Kirschbaum.2002b} davon auszugehen, dass eine rückläufige Entwicklung von Ausdrücken entscheidende Anhaltspunkte für den zur Diskussion stehenden Expressivitätsverlust bietet.

\begin{table}
\small
\begin{tabularx}{\textwidth}{l@{} YY YY YY}
\lsptoprule
{Intensivierer}  & \multicolumn{5}{c}{\textbf{Relative Frequenz}} & {Summe} \\
\cmidrule{2-6}
& 1950--1969 & Tendenz (in \%) & 1970--1989 & Tendenz (in \%) & 1990--2017 & \\
\midrule \\
\textit{krass} & -- & -- & 0,10 & +6\,910,00 & 7,01 & {7,10} \\
\textit{hammer} & -- & -- & 0,12 & +6\,216,67 & 7,58 & {7,69} \\
\textit{mega} & -- & -- & 0,24 & +4\,562,50 & 11,19 & {11,43} \\
\textit{ober} & \cellcolor[gray]{0.75} 3,01 & \cellcolor[gray]{0.75} - 5,98 & \cellcolor[gray]{0.75} 2,83 & +412,72 & 14,69 & {20,53} \\
\textit{derb}  & \cellcolor[gray]{0.75} 11,71 & \cellcolor[gray]{0.75} - 48,76  & \cellcolor[gray]{0.75} 6,00 & \cellcolor[gray]{0.75} - 27,50 & \cellcolor[gray]{0.75} 4,35 & {22,07} \\
\textit{voll} & 1,68 & +87,50 & 3,15 & +910,16 & 31,82 & {36,65} \\
\textit{scheiß} & 1,18 & +1\,061,86 & 13,71 & +97,52 & 27,08 & {41,96} \\
\textit{fürchterlich} & 6,27 & +120,89 & 13,85 & +68,66 & 23,36 & {43,49} \\
\textit{sau} & 10,78 & +24,40 & 13,41 & +170,10 & 36,22 & {60,41}  \\
\midrule
\textit{wahnsinnig} & 6,90 & +278,99 & 26,15 & +299,73 & 104,53 & {137,58} \\
\textit{unheimlich} & 25,69 & +150,02 & 64,23 & +12,05 & 71,97 & {161,89} \\
\textit{furchtbar} & 35,73 & +25,92 & 44,99 & +114,25 & 96,39 & {177,10} \\
\textit{arg} & 38,21 & +74,95 & 66,85 & +20,33 & 80,44 & {185,50} \\
\textit{schrecklich} & 40,65 & +39,11 & 56,55 & +57,49 & 89,06 & {186,26} \\
\textit{verdammt} & 39,94 & +28,34 & 51,26 & +167,69 & 137,22 & {228,42} \\
\textit{super} & 56,01 & +118,03 & 122,12 & +74,53 & 213,13 & {391,26}\\
\midrule
In Summe & 277,76 & +74,81 & 485,56 & +96,86 & 956,04    	& 1\,719,36\\
\lspbottomrule
\end{tabularx}
\caption[Kumulierte Häufigkeitsverteilungen der expressiven Intensivierer im Zeitraum von 1950 bis 2017]{Kumulierte Häufigkeitsverteilungen der expressiven Intensivierer im Zeitraum von 1950 bis 2017 (in zunehmender Frequenz).}
\label{Fre4}
\end{table}

%%%%%%%%%%%%%%%%%

\noindent Bei der Betrachtung von Tabelle \ref{Fre4} stellt sich heraus, dass beinahe alle der untersuchten Intensivierer trotz möglicher Schwankungen im Korpus häufiger im Gebrauch sind als zu Beginn des Untersuchungszeitraums. So sind einzig zwei Ausnahmen auszumachen, deren Frequenz prozentual betrachtet (zumindest temporär) schwindet. Dabei handelt es sich einerseits um \textit{derb}, bei dem quantitativ sowohl vom ersten zum zweiten Intervall als auch vom zweiten zum dritten Intervall eine starke Abnahme zu verzeichnen ist. Bei einem Vergleich des ersten Zeitraums mit dem dritten ist bei \textit{derb} insgesamt eine negative Differenz von -\,62,85\,\% zu konstatieren. Andererseits fällt ebenso der Gebrauch von \textit{ober} zeitweise marginal ab: Bei diesem Ausdruck ist eine leicht negative Differenz von -\,5,98\,\% vom ersten zum zweiten Zeitabschnitt zu erkennen. Diese Einbuße gleicht sich vom zweiten zum dritten Intervall allerdings insofern wieder aus, als die Frequenz von \textit{ober} in den beiden Zeitschriften verstärkt zunimmt. Alles in allem steigt die kumulierte Häufigkeit aller Ausdrücke vom ersten zum dritten Zeitabschnitt um 244,13\,\% an, was im Einklang mit den in \figref{Dia-rest} dargestellten diachronen Häufigkeitsverteilungen steht. Demnach haben die expressiven Intensivierer über den Untersuchungszeitraum hinweg stark an Frequenz dazugewonnen.

Als Nächstes wird eine paarweise Begutachtung der Zeitabschnitte vorgenommen. Dadurch kann überprüft werden, wie sich die expressiven Intensivierer relativ zu den Intervallen verhalten und welche Entwicklungstendenzen daraus abgeleitet werden können. Begonnen wird mit einer Gegenüberstellung des ersten Zeitraums (d.\,h. von 1950 bis 1969) mit dem zweiten Zeitraum (d.\,h. von 1970 bis 1989), woraufhin der zweite Zeitraum mit dem dritten Zeitraum (d.\,h. von 1990 bis 2017) kontrastiert wird. Im Anschluss daran wird der Einteilung der beiden Frequenzgruppen genauer auf den Grund gegangen. Vor diesem Hintergrund ergibt sich die Frage, ob die untersuchten Ausdrücke in allen drei Zeitabschnitten derselben Frequenzgruppe angehören oder von der einen Frequenzgruppe in die andere übertreten. Um diese zu beantworten, sind die Ausdrücke gemäß dieser Einteilung in Tabelle \ref{tab:Freqgruppe} aufgeführt. Darin wird zunächst jeder Zeitabschnitt separat berücksichtigt (d.\,h. 1950 bis 1969, 1970 bis 1989 sowie 1990 bis 2017), woraus für jeden Ausdruck letztlich eine Gesamttendenz abgeleitet wird.

\subsubsubsection{\textbf{1950--1969} vs. \textbf{1970--1989}}

Der markanteste Unterschied zwischen dem ersten und zweiten Zeitabschnitt besteht darin, dass im ersten noch nicht alle der untersuchten Intensivierer im Korpus ausgebildet sind. So lassen sich in diesem Intervall bloß 13 Ausdrücke nachweisen, tauchen \textit{hammer}, \textit{krass} und \textit{mega} doch erst im zweiten Intervall zum ersten Mal auf. Diese treten dann zunächst einmal in nur geringer Frequenz auf, was sich in Tabelle \ref{Fre4} naturgemäß in einer übermäßig stark ausgeprägten Anstiegskurve niederschlägt. Demnach handelt es sich bei \textit{hammer}, \textit{krass} und \textit{mega} um Ausdrücke, die erst im Laufe des Untersuchungszeitraums im Korpus neu dazukommen sind. Was die Mitglieder der beiden Frequenzgruppen anbelangt, gibt es mit Ausnahme der rezent etablierten Intensivierer keine bemerkenswerten Veränderungen. Lediglich \textit{wahnsinnig} kommt im zweiten Zeitabschnitt vergleichsweise vier Mal so häufig im Korpus vor, wodurch es von der Gruppe der niedrigfrequenten in die der hochfrequenten Ausdrücke übergetreten ist. Durch diese Entwicklung hat die \glqq{}schreckliche\grqq{} Intensivierergruppe in diesem Intervall zugleich ein ebenso negativ konnotiertes Mitglied dazugewonnen, was zu einem Ausbau ebendieser führt. Daneben ist auffällig, dass \textit{fürchterlich} seine Frequenz im zweiten Intervall nicht nur verdoppelt hat, sondern auch an die Spitze der niedrigfrequenten Gruppe gestiegen ist. Andere Intensivierer dieses Clusters weisen dagegen einen bedeutend geringeren Zuwachs auf, wie z.\,B. bei \textit{sau} oder \textit{voll} zu sehen ist. Auf der Spitzenposition der Gruppe der hochfrequenten Intensivierer steht sowohl im ersten wie auch im zweiten Intervall \textit{super}, das sein Vorkommen stark ausgebaut hat und bereits zu dieser Zeit im Korpus mit über 100 Belegen nachgewiesen werden kann. Darüber hinaus finden sich mit \textit{arg}, \textit{furchtbar}, \textit{schrecklich}, \textit{unheimlich} und \textit{verdammt} Ausdrücke mit klar negativer Quellsemantik, die im zweiten Zeitabschnitt, wie soeben erwähnt, durch \textit{wahnsinnig} ergänzt werden. Diese hochfrequenten Intensivierer befinden sich beim Übergang vom ersten zum zweiten Intervall quantitativ im Anstieg, wenn auch nicht dergestalt, dass sie an den Wert von \textit{super} heranreichen. Haben sich die kumulierten Frequenzen der hochfrequenten Intensivierer durch die stark positive Entwicklung von \textit{super} sogar verdoppelt, sind die der niedrigfrequenten vom ersten zum zweiten Intervall kaum merklich angestiegen. Bei der Betrachtung der Häufigkeiten aller Ausdrücke ist ein starker Anstieg von mehr als 70\,\% auszumachen, was ein Indiz für eine dem Anschein nach unproblematischere Verwendung von expressiven Intensivierern in den Zeitschriften zu sein scheint.

\subsubsubsection{\textbf{1970--1989} vs. \textbf{1990--2017}}

Bei der Gegenüberstellung des zweiten Intervalls mit dem dritten ist besonders auffällig, dass die rezent aufgetretenen Intensivierer \textit{krass}, \textit{hammer} und \textit{mega} im letzten Zeitabschnitt über eine höhere Frequenz als \textit{derb} verfügen. Letzterer ist der einzige Ausdruck, dessen Entwicklung über den Untersuchungszeitraum hinweg konstant rückläufig war. Zu beobachten ist außerdem, dass die Mitglieder der Frequenzgruppen in beiden Intervallen gleich geblieben sind und sich bloß die häufigkeitsbedingte Anordnung und dadurch auch die Subgruppierung im Dendrogramm geändert hat. So ist z.\,B. der einstige Spitzenreiter \textit{fürchterlich} trotz geringer Frequenzzunahme im dritten Zeitabschnitt von \textit{scheiß}, \textit{voll} und \textit{sau} überholt worden, deren Häufigkeiten sich in diesem Intervall zum Teil sehr stark ausgebaut haben. Die niedrigfrequente Gruppe zeigt im Hinblick auf die kumulierten Häufigkeiten eine positive Entwicklung beim Übergang vom zweiten zum dritten Intervall. Diese zeigt sich in der Form, dass die Spitzenreiter \textit{voll} und \textit{sau} über 30 Korpustreffer aufweisen, sodass sich der Frequenzwert von ca. 14 Belegen vom zweiten Zeitabschnitt im dritten mehr als verdoppelt hat. Auch in der Gruppe der hochfrequenten Intensivierer treten derart markante Veränderungen auf. In dieser steht \textit{super} nach wie vor unangefochten an erster Stelle, während sich seine einstige Trefferzahl im dritten Intervall sogar auf über 200 Belege erhöht hat. Zudem finden sich mit \textit{wahnsinnig} und \textit{verdammt} zwei weitere Intensivierer, die mit über 100 Treffern pro Million Textwörter im Korpus nachgewiesen werden können. Bemerkenswert ist, dass die Frequenz vom ersten bzw. zweiten zum dritten Intervall stark zugenommen hat. So hat sich \textit{wahnsinnig} vom einstigen niedrigfrequenten Intensivierer im ersten Zeitabschnitt zu einem der höchstfrequenten Intensivierer im dritten Zeitabschnitt entwickelt, wohingegen sich ebenso die Frequenz von \textit{verdammt} im dritten Intervall mehr als verdoppelt hat. Demzufolge ist auch in diesem Cluster ein übermäßig ausgeprägter Anstieg der Ausdrücke zu verzeichnen, was eine stark positive Entwicklung in den kumulierten Häufigkeiten aller Intensivierer nach sich zieht. Die Gesamtfrequenz der expressiven Intensivierer hat sich beim Übergang vom zweiten zum dritten Intervall ebenfalls nahezu dupliziert. Infolgedessen sind letztlich mehr als doppelt so viele Intensivierer im Zeitschriftenkorpus enthalten wie zu Beginn des Untersuchungszeitraums. Dies ist auch mithilfe des in \figref{Dia-rest} dargestellten Säulendiagramms zu beobachten (vgl. erster und letzter Balken).

Neben den kumulierten Häufigkeiten werden die expressiven Intensivierer, wie oben angekündigt, mit Blick auf die beiden Frequenzgruppen `Niedrigfrequent' und `Hochfrequent' begutachtet. Tabelle \ref{tab:Freqgruppe} gibt einen Überblick über die Mitglieder der Frequenzgruppen je Zeitabschnitt und leitet auf dieser Grundlage für jeden Intensivierer eine Gesamttendenz ab. Die Sortierung der Ausdrücke erfolgt in ansteigender Frequenz. In der Tabelle ist jeweils diejenige Spalte farblich hervorgehoben, bei der ein Wechsel der Frequenzgruppe zu verzeichnen ist.

\begin{table}
\begin{tabularx}{\textwidth}{lXXXrl}
\lsptoprule
\multirow{3}{*}{Intensivierer} &  \multicolumn{3}{c}{\textbf{Frequenzgruppe}}  & & \multirow{3}{*}{Klassifikation}    \\
 & 1950--1969 & 1970--1989 & 1990--2017 &  &  \\
\midrule
\textit{hammer}  &  {--}	& {niedrig} & {niedrig}	& \rdelim\}{9}{.3cm} & \\
\textit{krass}  &  {--}	& {niedrig} & {niedrig}	&  &  \\
\textit{mega}  &  {--}	& {niedrig} & {niedrig}	&  &  \\
\textit{derb}  &  {niedrig}	& {niedrig} & {niedrig}	&  &   \\
\textit{fürchterlich}  & {niedrig}	& {niedrig} & {niedrig}	&  &   { {{niedrig   (n = 9)}}}  \\
\textit{ober}  &  {niedrig}& {niedrig} & {niedrig}	&  &  \\
\textit{sau}  & {niedrig}	& {niedrig} & {niedrig}	&  &  \\
\textit{scheiß}  & {niedrig}	& {niedrig} & {niedrig}	&  &  \\
\textit{voll}  & {niedrig}	& {niedrig} & {niedrig}	&  &  \\
\midrule
\shadecell  \textit{wahnsinnig}  & \shadecell   niedrig &   \shadecell   hoch  &   \shadecell   hoch 	& \rdelim\}{7}{.3cm}  &  \\
\textit{arg}  &  {hoch}	& {hoch} & {hoch}	& &  \\
\textit{furchtbar}  &  {hoch}	& {hoch} & {hoch}	& & \\
\textit{schrecklich}  &  {hoch}	& {hoch} & {hoch}	& &   { {{hoch   (n = 7)}}}  \\
\textit{super}  &  {hoch}	& {hoch} & {hoch}	& &  \\
\textit{unheimlich}  &  {hoch}	& {hoch} & {hoch}	& &  \\
\textit{verdammt}  &  {hoch}	& {hoch} & {hoch}	& &  \\
\midrule\\
\lspbottomrule
\end{tabularx}
\caption[Einteilung der expressiven Intensivierer nach Frequenz im Zeitraum von 1950 bis 2017]{Einteilung der expressiven Intensivierer nach Frequenz im Zeitraum von 1950 bis 2017 (angeordnet nach Frequenzgruppe).}
\label{tab:Freqgruppe}
\end{table}

Aus Tabelle \ref{tab:Freqgruppe} ist abzulesen, dass die meisten expressiven Intensivierer ihre Frequenzgruppe über den gesamten Untersuchungszeitraum beibehalten. So zeichnet sich bei 15 Intensivierern dahingehend ein konstantes Muster ab, wobei niedrig- und hochfrequente Ausdrücke gleichermaßen davon betroffen sind. Eine Sonderstellung nimmt \textit{wahnsinnig} ein, das beim Wechsel vom ersten zum zweiten Zeitabschnitt von der niedrigfrequenten Gruppe in die Gruppe der hochfrequenten Intensivierer übergetreten ist. Weil der Ausdruck im Korpus auch im dritten Intervall unverändert hochfrequent ist und sich die im zweiten Intervall ermittelte Häufigkeit von 26,15 Treffer pro Million Textwörter im dritten Intervall auf 104,53 Belege vervierfacht hat (vgl. die Tabellen \ref{Fre2} und \ref{Fre3.1}), ist seine Gesamttendenz ebenso als hoch einzuschätzen.

Zusammengenommen umfassen die Gruppen ingesamt neun niedrigfrequente und sieben hochfrequente Intensivierer. Lassen sich erstere in Summe mit 258,24 Treffern pro Million Textwörter nachweisen ($\hateq$ 15,02\,\%), sind letztere im Zeitschriftenkorpus mit 1461,12 Belegen vertreten ($\hateq$ 84,98\,\%).

\section{Diskussion} 		\label{DisKor}
\largerpage
Nachfolgend werden die im vorangegangenen Abschnitt dargelegten Ergebnisse der Korpusuntersuchung besprochen. Diese wurde mit dem Ziel durchgeführt, synchrone und diachrone Häufigkeitsverteilungen ausgewählter expressiver Intensivierer über einen umfassenden Zeitraum zu erheben, um einigen in der Literatur vertretenen Annahmen zum ersten Mal empirisch nachzugehen. So ist im Sinne der Hypothese H1-Freq nach \citet[2]{Hauschild.1899} davon auszugehen, dass Intensivierer infolge von inflationärem Gebrauch ihre Expressivität verlieren und in der Konsequenz seltener verwendet werden. Nach \citet[59]{Tobler.1868} spielt dabei auch der Faktor Zeit eine Rolle, der in der Hypothese H1-Temp berücksichtigt wurde. Um den Ausfall abgenutzter Intensivierer zu kompensieren, bedarf es \citet[274]{Baumgarten.1908} zufolge permanent neuer expressiver Ausdrücke. Diese Aussage wurde im Rahmen der Hypothese H2 auf ihre Gültigkeit hin überprüft.

Nachdem im Folgenden zunächst auf den deskriptiven Intensivierer \textit{sehr} eingegangen wird, folgt daraufhin eine Begutachtung der expressiven Intensivierer unter Bezugnahme der Ergebnisse der hierarchischen Clusteranalyse. Abschließend werden die gewonnenen Resultate zusammengefasst und in Hinsicht auf die soeben umrissenen Hypothesen evaluiert.

\subsection{Deskriptiver Intensivierer}

Anhand der Ergebnisse der Korpusanalyse lässt sich die longitudinale Entwicklung von \textit{sehr} als eine Delle mit fallendem Beginn und steigendem Ende beschreiben (vgl. \figref{Dia-sehr}), was einem diachronen Jo-Jo-Effekt entspricht. Der Tiefpunkt, der sich beim deskriptiven Intensivierer im Zeitraum von 1970 bis 1990 detektieren lässt, ist also überwunden. Durch Extrapolation ist anzunehmen, dass die Frequenzzunahme weiter andauert und das Gebrauchsmaximum noch nicht erreicht ist -- vorausgesetzt, die zuletzt beobachteten Häufigkeiten von \textit{sehr} bleiben in den darauffolgenden Jahren erhalten. Mit 20\,278,12 Belegen pro Million Textwörter ist \textit{sehr}, wie prognostiziert, der mit Abstand am meisten verwendete Intensivierer im Korpus. Dieses Ergebnis deckt sich mit der Beobachtung, die \citet[266]{Ito.2003} im Rahmen ihrer Korpusuntersuchung in Bezug auf das englische Pendant \textit{very} machten. Auch legt es die Vermutung nahe, dass \textit{sehr} unangefochten der häufigste Intensivierer im Deutschen ist. Somit besteht offensichtlich keine Gefahr eines dauerhaften Rückgangs oder gar einer Verdrängung durch andere Intensivierer. Angesichts der anfangs zu verzeichnenden Frequenzabnahme von \textit{sehr} lässt sich einerseits spekulieren, dass es sich dabei um ein Artefakt der Daten handelt, das nicht signifikant ist und aus linguistischer Sicht keine Relevanz hat. Andererseits können ebenfalls die gesellschaftlichen Rahmenbedingungen als Erklärung in Betracht gezogen werden. So ist z.\,B. denkbar, dass der Schreibstil in den Zeitschriften in der Nachkriegszeit nüchterner ausfiel, die Journalist:innen neutralere Berichterstattung leisteten und der Ausdruck einer allgemeinen Tendenz zur Vermeidung von Intensivierern unterlag. Dementsprechend könnte die Ursache für den zeitweisen Schwund extralinguistisch durch soziokulturelle Faktoren bedingt sein. Auch wenn bezüglich des Rückgangs an dieser Stelle lediglich Spekulationen möglich sind, deutet der Anstieg, der sich ab den Neunzigerjahren abzeichnet, auf eine weiterhin starke Verwendung ohne Abnutzungserscheinungen hin. Diese ist abgesehen vom hohen Lexikalisierungsgrad vor allem der skalaren Bedeutung geschuldet, die \textit{sehr} zum klassischen Intensivierer des Standarddeutschen macht. Im Hinblick auf die eher konservativen Textsorten, aus denen sich das Untersuchungskorpus zusammensetzt, bietet \textit{sehr} die neutralste Möglichkeit zur Skalierung, die auch in Printmedien angemessen ist. So handelt es sich bei Zeitungen nach \citet[]{Koch.1996,Koch.2007} um eine konzeptionell schriftliche Kommunikationsform, die deutlich stärker zu einem deskriptiven als expressiven Stil neigt. Dies konnte durch die Untersuchung gezeigt werden. Dagegen gilt die expressive Intensivierung gemeinhin als mündlichkeitsnahes Phänomen, das, ähnlich wie andere Formen von Expressivität (z.\,B. Modalpartikeln), in der Schriftsprache eher zu vermeiden ist \citep[vgl.][9]{Weydt.1969}. Obwohl der Gebrauch von \textit{sehr} nach einer temporären Einbuße aus den genannten Gründen (d.\,h. stilistisch wenig stigmatisierte Verwendung durch Delexikalisierung, hoher Lexikalisierungsgrad und fehlende Expressivität) im Korpus zwar wieder zugenommen hat, ist dennoch vorstellbar, dass es in der mündlichen Kommunikation nicht ausdrucksstark genug ist. Infolgedessen ist gemäß den in Abschnitt \ref{Auslöser} dargelegten Expressivitätsmaximen nach \citet[]{Kirschbaum.2002a} in mündlichkeitsnahen Sprachdomänen zu erwarten, dass expressive Intensivierer gegenüber deskriptiven Intensivierern präferiert verwendet werden, wenn zusätzlich zur Graduierung auch subjektive Einstellungen bekundet werden sollen -- jedenfalls sofern die hierzu in der Literatur vorherrschende Annahme (vgl. bspw. \citealt[315]{Lang.1983}; \citealt[57]{Ortner.2014}) korrekt ist. Die noch bestehenden Unklarheiten und angeführten Mutmaßungen in Bezug auf die diachrone Entwicklung von \textit{sehr} können als theoretische Grundlage für weiterführende korpuslinguistische oder experimentelle Untersuchungen dienen.

\subsection{Expressive Intensivierer}

Im Anschluss werden die Ergebnisse der expressiven Intensivierer besprochen. Die Diskussion erfolgt zum einen auf Basis der festgestellten semantischen Ähnlichkeiten und zum anderen auf Basis der kumulierten Häufigkeiten unter Bezugnahme der Ergebnisse der hierarchischen Clusteranalyse.

Das erste augenscheinliche Resultat hinsichtlich der expressiven Intensivierer ist mithilfe des in \figref{Dia-rest} dargestellten Säulendiagramms zu erkennen. Dem darin veranschaulichten diachronen Verlauf folgend haben die expressiven Intensivierer stark an Frequenz dazugewonnen. So sind die in den ersten Intervallen verzeichneten Häufigkeiten bis zu den letzten linear und beinahe auf das Dreifache angestiegen. Wie an entsprechender Stelle ausgeführt, stellt dies insbesondere vor dem Hintergrund des zurate gezogenen Korpus ein unerwartetes Ergebnis dar. Somit hat nicht nur der Gebrauch des deskriptiven Intensivierers \textit{sehr} seit den Neunzigerjahren konstant zugenommen, sondern auch die expressiven Intensivierer zeigen ähnliche Entwicklungstendenzen an. Dies gibt erste Hinweise auf eine mögliche stilistische Veränderung der in den Zeitschriften verwendeten Schreibweise. Ebenso ist nicht ausgeschlossen, dass die diachronen Entwicklungen durch eine stärkere Gewichtung zugunsten meinungsäußernder Texte bzw. konzeptionell mündlich orientierter Textsorten wie Interviews, Leserbriefen oder Kommentaren begründet sind. Diese Annahme lässt sich allerdings aufgrund unzureichender Metadaten der Korpusdokumentation nicht belegen. Ungeachtet der Gründe ist die Frequenzzunahme unstrittig. Als Ergebnis kann folglich festgehalten werden, dass die Journalist:innen nach und nach vermehrt auf mündlichkeitsnahe Ausdrücke zurückgegriffen haben, was einen unkritischeren Gebrauch expressiver Intensivierer im Korpus nahelegt.

Zur Prüfung der Hypothesen H1-Temp und H1-Freq wurde der Untersuchungszeitraum in drei Abschnitte gegliedert, in denen die Häufigkeiten systematisch inspiziert wurden. Hier zeichneten sich recht früh erste Entwicklungstendenzen ab, die im Folgenden betrachtet werden. Begonnen wird mit den frequenten Intensivierern, allen voran \textit{super}. Dieser Ausdruck war von Beginn an am häufigsten im Korpus vertreten. Bis zum dritten Intervall hat sich seine Frequenz mit deutlicher Distanz zu den anderen Ausdrücken ausgebaut. Bei \textit{super} handelt es sich um einen Intensivierer lateinischen Ursprungs, der anzeigt, dass eine kontextuell determinierte Schwelle überschritten wird, sodass die beschriebene Eigenschaft eine \glqq{}Abweichung von der Normalität\grqq{} (\citealt[80]{Ruf.1996}) darstelle. Der Umstand, dass ausgerechnet der frequenteste Intensivierer ein Lehnwort ist, entkräftet überdies zu einem gewissen Grad die eingangs angeführte Hypothese von \citet[323]{Calpestrati.2017}, nach der deutsche Muttersprachler:innen in der Schriftsprache heimische Ausdrücke gegenüber entlehnten bevorzugen. Wenngleich dies auf die restlichen hochfrequenten Intensivierer durchaus zutrifft, stehen die Frequenzen von \textit{super} jedoch im Widerspruch dazu. In Bezug auf dessen Etymologie nach \citet[1398]{Pfeifer.1995} ist festzustellen, dass es sich bereits im 16. Jahrhundert als Intensivierer etabliert hat. Daher ist davon auszugehen, dass der Ausdruck einer der ältesten der untersuchten Intensivierer ist,\footnote{Ein weiteres Indiz für die lange Existenz des intensivierenden \textit{super} findet sich in \citet[109]{Tobler.1868}, der \textit{superklug} neben \textit{hyperkritik}, \textit{extrafein} und \textit{ultraradical} als Beispiel für eine \glqq{}moderne bildung\grqq{} nennt.} was als Erklärung für die verhältnismäßig hohe Korpusfrequenz in Erwägung gezogen werden kann. Interessanterweise deckt sich dieses Ergebnis mit den Resultaten, die \citet[345]{Ruf.1996} und \citet[104]{Androutsopoulos.1998} im Rahmen ihrer Korpusstudien erzielten: Auch bei ihnen trat \textit{super} mit Abstand am häufigsten im Korpus auf. \citet[345]{Ruf.1996} verweist in diesem Zusammenhang auf \citet[]{Wunderlich.1983}, der \textit{super} eine \glqq{}große Zukunft\grqq{} prophezeite. Dass der Anstieg bislang nicht abgeschlossen, das Frequenzmaximum noch nicht erreicht und der Intensivierer auch gegenwärtig von großer Relevanz ist, demonstrieren die Ergebnisse meiner Untersuchung ganz klar. Doch wie ist es um die Expressivität von \textit{super} bestellt? Dies lässt sich auf Basis der quantitativen Korpusdaten nicht abschätzen, sondern erfordert eine experimentelle Überprüfung. Eine solche vertiefende Einordnung wird in der in Kapitel \ref{Experiment0} vorgestellten experimentellen Untersuchung vorgenommen.

Ein gleichfalls markantes Ergebnis stellt die weitere Zusammensetzung des Clusters der hochfrequenten Ausdrücke dar, das neben \textit{super} schon von Beginn des Untersuchungszeitraums an \textit{unheimlich}, \textit{furchtbar}, \textit{arg}, \textit{verdammt} und \textit{schrecklich} beinhaltete. Zeigte dieses Cluster bereits im ersten Intervall eine Tendenz für Intensivierer mit klar negativer Semantik, wurde es ab dem zweiten Intervall zusätzlich durch \textit{wahnsinnig} unterstützt; einen Ausdruck, der ebenso einen negativen mentalen Zustand denotiert. Ähnlich wie \textit{super} bewegen sich diese Ausdrücke in einem hohen Frequenzbereich. Bei ihnen handelt es sich ausschließlich um transparente denominale bzw. deadjektivische Derivationen. Obwohl deren Semantik zwar erkennbar ist, ist sie in intensivierender Funktion der Einschätzung von \citet[15f.]{Nouwen.2020} zufolge jedoch gänzlich verblasst. Auch bei diesen Ausdrücken lässt sich die intensivierende Funktion mittels \mbox{\citet[385, 525, 1499]{Pfeifer.1995}} zeitlich früh verorten, d.\,h. zwischen dem 16. Jahrhundert (\textit{furchtbar}) und dem 19. Jahrhundert (\textit{unheimlich} und \textit{verdammt}). Demnach ist es naheliegend, dass die Intensivierer inzwischen über einen hohen Lexikalisierungsgrad verfügen und eine eher skalare Bedeutung angenommen haben. Dieser Punkt wird weiter unten unter dem Blickwinkel der Definition von Expressivitätsverlust erneut aufgegriffen und vertieft. Bei diesen Ausdrücken stellt sich also ebenfalls die Frage nach dem tatsächlichen Expressivitätsgrad. Vor dem Hintergrund der Kriterien Erstvorkommen (vgl. \mbox{H1-Temp}) und Frequenz (vgl. H1-Freq) ist anzunehmen, dass die hochfrequenten Intensivierer prinzipiell weniger expressiv sind als die niedrigfrequenten. Auch darüber kann die nachfolgend dargestellte experimentelle Untersuchung Aufschluss geben.

Wesentlich variabler ist dagegen die Zusammensetzung der Gruppe der niedrigfrequenten Intensivierer. Hier findet sich u.\,a. der einzige Ausdruck, der über die Zeit als Intensivierer sukzessive seltener gebraucht wurde, nämlich \textit{derb}. Dessen Verwendung ist seit dem Frequenzhöhepunkt Anfang der Sechzigerjahre weitgehend kontinuierlich zurückgegangen, was schließlich dazu führte, dass der Intensivierer im Korpus nahezu ausgestorben ist. Folglich stellt die beständige diachrone Frequenzeinbuße von \textit{derb} als Einziges ein hinreichendes Indiz für einen Verlust von Expressivität dar. Ob sich die Frequenzabnahme darauf zurückführen lässt, dass der Ausdruck infolge früherer häufiger oder persistenter Verwendung nach und nach an Expressivität verloren hat, oder ob er durch andere Intensivierer verdrängt wurde, weil er schlichtweg unbeliebter war, lässt sich korpuslinguistisch nicht überprüfen. In beiden Fällen könnte der für \textit{derb} nachteilige Status den Wandel zum Vorteil anderer Ausdrücke beschleunigt haben. Gemessen an den IpMW ist \textit{derb} neben dem vergleichsweise neuen Ausdruck \textit{krass} im letzten Zeitabschnitt (d.\,h. von 1990 bis 2017) der am seltensten verwendete Intensivierer. Auch die Betrachtung des ersten Intervalls (d.\,h. von 1950 bis 1969) bringt keine Indizien für einen hochfrequenten Gebrauch hervor. Damit bleibt die Frage bestehen, ob \textit{derb} in der Zeit vor 1950 häufiger (und unter Umständen zu häufig) verwendet wurde (vgl. H1-Freq). Diese muss in Anbetracht des beschränkten Untersuchungszeitraums allerdings unbeantwortet bleiben. Ähnliches trifft auf die Persistenz zu, die sich gleichermaßen als nachteilig für die Expressivität erweisen kann \mbox{(vgl. H1-Temp)}. Ungeachtet der tatsächlichen Gründe ist \textit{derb} auf den ersten Blick der einzige Intensivierer, dessen fortschreitender quantitativer Rückgang im Zeitschriftenkorpus zweifelsfrei auf einen Expressivitätsverlust hindeutet. Außerdem befinden sich in der Gruppe der niedrigfrequenten Ausdrücke diejenigen Intensivierer, die sich erst im Laufe der Zeit im Korpus neu etabliert haben und somit Evidenz für die Gültigkeit der Hypothese H2 nach \citet[274]{Baumgarten.1908} liefern, nämlich \textit{krass}, \textit{hammer} und \textit{mega}. Derzufolge gibt es einen ständigen Bedarf an neuen Intensivierern, um den Ausfall altgedienter Ausdrücke zu kompensieren. Darin sei zugleich die Hauptursache für die Wandelbarkeit des Inventars zu sehen, die dazu führt, dass Intensivierer in der Literatur als \glqq{}Modewörter\grqq{} (\citealt[267]{Becher.1907}) bezeichnet werden. Wie eingangs beschrieben, sollte sich das Zutreffen von H2 im Korpus durch die Herausbildung neuer Intensivierer ausdrücken. Dies legt die Entwicklung von \textit{krass}, \textit{hammer} und \textit{mega} nahe. Semantisch betrachtet entspricht letzterer Ausdruck der deutschen Präposition \textit{ober}, das ebenso in dieser Gruppe zu finden ist. Während \textit{mega} über das Griechische entlehnt wurde und nach \mbox{\citet[855]{Pfeifer.1995}} so viel wie `groß, hoch, lang, weit' bedeutet, trägt auch \textit{ober} eine dimensionale Bedeutung, die \mbox{\citet[940]{Pfeifer.1995}} mit `oben befindlich, höher gelegen, auf höherer Stufe stehend, übergeordnet' angibt. Daneben können \textit{hammer} und \textit{krass} im weiteren Sinne ebenso als verwandt angesehen werden: Der Hammer stellt als robustes Schlagwerkzeug einen starken, plumpen Gegenstand dar, der aufgrund seiner massiven physischen Kraft dazu dient, grobe Arbeiten zu verrichten. Ebendiese Adjektive, d.\,h. `dick, grob, derb, plump', bringt \citet[727]{Pfeifer.1995} mit der etymologischen Bedeutung von \textit{krass} in Verbindung, weshalb hier ebenfalls eine semantische Relation zu sehen ist. Bei \textit{fürchterlich}, \textit{scheiß} und \textit{sau} ist auffällig, dass es sich um Ausdrücke handelt, die sich ursprünglich von den polysemen, überwiegend negativ konnotierten Nomen abgeleitet haben: von der \textit{Sau} als weiblichem Schwein, der emotionalen Reaktion \textit{Furcht} sowie dem Fäkalwort \textit{Scheiße}. Insbesondere bei letzterem könnte die semantische Transparenz und die damit verbundene vulgäre Konnotation ausschlaggebend dafür sein, dass der Ausdruck im Zeitschriftenkorpus nicht häufiger vertreten ist. Ferner wird auch \textit{sau} in schriftsprachlichen Kontexten offenbar seit jeher als unangemessen betrachtet, tritt es nach \mbox{\citet[122]{Tobler.1868}} doch \glqq{}nur in der gemeinern sprache der mundarten\grqq{} auf. Auch existiert der Ausdruck seit dem Mittelhochdeutschen als Schimpfwort: \textit{Du blöde Sau} \mbox{\citep[vgl.][1168]{Pfeifer.1995}}.\footnote{Trotzdem berichtet \citet[270]{Becher.1907}, dass bei \textit{Schwein} und \textit{Sau} ehedem die Symbolik im Vordergrund gestanden habe, wurde das Tier in der Geschichte doch als heilig verehrt. Heute spielt hingegen wohl eher die damit verbundene negative Konnotation eine übergeordnete Rolle, durch die die Ausdrücke nahezu vorherbestimmt für die \mbox{Intensiviererrolle sind}.} Bei der Interpretation von Ausdrücke tierischen Ursprungs sieht \mbox{\citet[]{Claudi.2006}} primär die Gegensatzrelation Mensch vs. Tier bzw. Zivilisation vs. Natur vordergründig:

\begin{singlespace*}
\begin{quotation}
\noindent The intensifiers from the animal domain probably derived their intensifying power from a contrast humanity vs. non-humanity, or culture/civilisation vs. nature, where non-humanity or nature stands for everything that is uncontrolled, untamed, and instinctual. \hfill (\citealt[360]{Claudi.2006}) \\ \end{quotation}
\end{singlespace*}

Gleichfalls finden sich bei \textit{fürchterlich} und \textit{sau} in \mbox{\citet[386, 1168]{Pfeifer.1995}} Hinweise darauf, dass die Ausdrücke erstmals im 17. bzw. 19. Jahrhundert intensivierend aufgetreten sind. Das Erstvorkommen der anderen Intensivierer ist dagegen nicht belegt. Auf den ersten Blick ist es überraschend, dass \textit{fürchterlich} im Korpus nicht frequenter ist, steht es den Mitgliedern der \glqq{}schrecklichen\grqq{} Intensivierergruppe zumindest semantisch in Nichts nach. Nichtsdestotrotz handelt es sich bei vielen Intensivierern um morphologisch und phonologisch wenig komplexe Ausdrücke, die sich im Affekt schnell, unkompliziert und dennoch mit dem intendierten kommunikativ-pragmatischen Effekt artikulieren lassen. Vergleicht man bspw. die neuen Intensivierer \textit{krass} oder \textit{mega} mit den älteren wie \textit{unheimlich} oder \textit{wahnsinnig}, so zeichnet sich ein deutlicher Trend zur Kürze ab. Wie bereits erwähnt, sind die hohen Frequenzen der älteren Intensivierer vermutlich in erster Linie auf einen hohen Lexikalisierungsgrad zurückzuführen. Im Falle von \textit{fürchterlich} lässt sich an dieser Stelle mutmaßen, dass es zugunsten von \textit{furchtbar} als morphologisch und phonologisch einfacherer Variante (da weniger Silben und kein Umlaut) seltener gebraucht wurde. Zu guter Letzt fehlt noch \textit{scheiß}, bei dem 82,40\,\% der Treffer die weitgehend konventionalisierte Komposition \textit{scheißegal} (oder auch die Berliner Variante mit Phonemverschiebung von [ɡ] zu [j], d.\,h. \textit{scheißejal}) ausmachen, wohingegen lediglich 17,60\,\% auf andere Verbindungen zukommen. Daher ist zu erwägen, dass der Intensivierer \textit{scheiß} auf einige wenige Adjektive limitiert und gegebenenfalls nicht bloß im Korpus wenig produktiv ist. Bezüglich der Produktivität von niedrigfrequenten Ausdrücken erklärt \citet[72]{Kempf.2016}, dass eine \glqq{}sehr niedrige Tokenfrequenz einer Bildung darauf hindeuten [kann], dass es sich um eine Neubildung handelt.\grqq{} Meiner Ansicht nach ist dies bei \textit{scheiß} allerdings aus verschiedenen Gründen ausgeschlossen: Zum einen finden sich in meinem Korpus mit \textit{scheißfreundlich} und \textit{scheißnormal} schon im Zeitraum von 1958 bis 1960 erste Belege für das intensivierende \textit{scheiß}. Zum anderen hat sich sein Vorkommen über die Zeit allenfalls marginal verändert. Demnach ist der verhältnismäßig niedrige Gebrauch anderweitig begründet. So ist anzunehmen, dass die negative Konnotation des Ausdrucks in der Schriftsprache weniger angemessen ist als im mündlichen Sprachgebrauch, was dazu führt, dass \textit{scheiß} im Vergleich zu weniger offensiven Intensivierern seltener verwendet wird. Dass offensive, negativ konnotierte Ausdrücke wie \textit{scheiß} oder \textit{arsch} über eine derart geringe Korpusfrequenz verfügen, legt obendrein die Existenz eines semantischen Filters nahe. Dieser scheint in der Schriftsprache dafür zu sorgen, dass expressive Intensivierer mit negativem Herkunftsbereich und ebensolcher Konnotation schlechter Einzug in dieses Medium halten als Intensivierer, die opak oder abgenutzter, somit textsortenadäquater sind.

Auf Basis der Ergebnisse lässt sich zusammenfassen, dass die Häufigkeiten anscheinend zumindest partiell durch die Semantik der Ausdrücke zu erklären sind. Dieser Einfluss kann sowohl von Vorteil als auch von Nachteil sein: Während die semantische Transparenz die hohe Frequenz der \mbox{Mitglieder} der \glqq{}schrecklichen\grqq{} Intensivierergruppe womöglich begünstigt hat, hat sie sich bei anderen Intensivierern dagegen negativ ausgewirkt. Dies ist z.\,B. bei \textit{sau} und \textit{scheiß} zu beobachten, die sich im Korpus in nur geringer Zahl nachweisen lassen. Bei ihnen hat die Transparenz und die damit verbundene negative Konnotation den Gebrauch der Ausdrücke im Korpus offenbar behindert und dazu geführt, dass andere, weniger offensive Intensivierer bevorzugt verwendet wurden. So konstatiert bspw. \citet[311]{Calpestrati.2017}, dass der semantischen Konnotation eines Intensivierers eine tragende Rolle zukommt. Dies bestätigen auch die Ergebnisse meiner Korpusanalyse. Dennoch hebt \citet[229]{Pittner.1989} die Möglichkeit hervor, dass sich die mit einem Ausdruck verbundene Konnotation je nach Sprecher und Äußerungskontext verändern könne. So sei \textit{sau} für Erwachsene grundsätzlich stärker negativ konnotiert, wohingegen es für Jugendliche vielmehr mit etwas Positivem verbunden werde (\textit{sauteuer}, \textit{saudumm} vs. \textit{saugeil}, \textit{saustark}) \citep[vgl.][229]{Pittner.1989}. Ob diese dem Anschein nach introspektive Einschätzung (auch aus heutiger Sicht) der Realität entspricht, bedarf empirischer Überprüfung.

Doch was bedeuten die erbrachten Resultate für die zu testenden Hypothesen von \citet[2]{Hauschild.1899} und \citet[274]{Baumgarten.1908}, nach denen eine persistente und frequente Verwendung einen Verlust der Expressivität nach sich zieht (vgl. H1), woraus sich schließlich ein ständiger Bedarf an neuen Ausdrücken ergibt (vgl. H2)? Wie sich anhand der etwa 70-jährigen Entwicklung ablesen lässt, sind derzeit 15 der 16 ausgewerteten Intensivierer quantitativ im Anstieg. Dieser Anstieg bietet auf den ersten Blick zunächst einmal keine direkten Anhaltspunkte für einen Expressivitätsverlust. So ist vor dem Hintergrund der Expressivitätsmaximen nach \mbox{\citet[]{Kirschbaum.2002b}} anzunehmen, dass ein Ausdruck, der an Expressivität eingebüßt hat, bei Sprechern weniger beliebt ist, von anderen verdrängt und letztlich seltener verwendet wird. Der einzige Intensivierer, bei dem dies dem diachronen Verlauf nach zutreffend ist, ist \textit{derb}. Obwohl die genauen Gründe für den quantitativen Rückgang von \textit{derb} nicht zweifelsfrei benannt werden können, lassen die Ergebnisse doch insofern Raum für Spekulation, als dem zur Diskussion stehenden Expressivitätsverlust zumindest eine tragende Rolle als treibende Kraft zugesprochen werden kann. Nichtsdestotrotz besteht neben der beschriebenen Definition von Expressivitätsverlust als Frequenzrückgang ebenso die Möglichkeit, dass er sich davon abweichend -- vielleicht sogar in Form eines Frequenzanstiegs -- niederschlägt: Erweist sich das soeben beschriebene Szenario wahrscheinlich im Kontext gesprochener Sprache als richtig, äußert sich ein Verlust der Expressivität in der Schriftsprache unter Umständen anders, wenn nicht sogar gegenläufig. Demzufolge unterscheiden sich die Auswirkungen des Expressivitätsverlusts eventuell je nach Medium der Realisierung im Sinne von \mbox{\citet[]{Koch.1996,Koch.2007}}, d.\,h. gesprochene vs. geschriebene Sprache. Im Hinblick auf das Zeitschriftenkorpus ist davon auszugehen, dass ein zu hoher Expressivitätsgrad nicht vereinbar mit den Eigenschaften der darin enthaltenen neutralen Textsorten ist. Besser eignen sich hingegen opake, weniger expressive oder pragmatisch abgeschliffene Intensivierer, was sich möglicherweise in einer Zunahme deren Gebrauchsfrequenz äußert. Eine solche Entwicklung würde bspw. ebenso die stetige Frequenzzunahme der Gruppe der \glqq{}schrecklichen\grqq{} Intensivierer wie \textit{verdammt} oder \textit{wahnsinnig} erklären. Letzterer wechselte sogar als einziger Ausdruck seine Frequenzgruppe (vgl. Tabelle \ref{tab:Freqgruppe}): Gehörte er im ersten Intervall noch zu den niedrigfrequenten Intensivierern, war er ab dem zweiten als hochfrequent einzuschätzen und entwickelte sich schließlich zu einem der häufigsten Ausdrücke im Korpus. Im Falle der hochfrequenten Intensivierer könnte der konstante quantitative Anstieg im Korpus also nicht bloß auf einen hohen Lexikalisierungsgrad, sondern auch auf einen fortgeschrittenen Expressivitätsverlust hindeuten. Auf dieser Grundlage ist vorstellbar, dass die hochfrequenten Ausdrücke sowohl an Expressivität als auch an lexikalischen Eigenschaften eingebüßt haben und in der Folge stilistisch weniger mit der neutralen Zeitungssprache in Konflikt geraten als offenkundig expressive(re) Ausdrücke wie z.\,B. \textit{scheiß} oder das aufgrund einer zu geringen Belegzahl aussortierte \textit{arsch}. Demnach resultiert der Expressivitätsverlust im schriftlichen Sprachgebrauch gegebenenfalls in einer besseren Eignung, einer größeren Flexibilität in der Verwendung sowie einer entsprechend höheren Frequenz in den Zeitschriften. Mit anderen Worten: Je niedriger der Expressivitätsgrad eines Intensivierers ist, in desto mehr schriftsprachlichen Kontexten ist er vermutlich anwendbar. Diese Annahme steht überdies im Einklang mit der Aussage von \mbox{\citet[420]{Breindl.2007}}, nach der sich der Anwendungsbereich eines Ausdrucks mit fortschreitender Delexikalisierung ausdehnt, was gleichzeitig in einem Anstieg dessen Gebrauchs resultiert. Diese Abweichung von der nach \citet[]{Kirschbaum.2002b} zu erwartenden Entwicklung legt medienspezifische Unterschiede nahe: Während sich der Expressivitätsverlust in der gesprochenen Sprache in Form einer rückläufigen Entwicklung bemerkbar macht, bewirkt er in der geschriebenen Sprache stattdessen potenziell einen Anstieg der Frequenz -- vor allem in stilistisch konservativen Textsorten. Wenn die Verwendung von Intensivierern in schriftlichkeitsnahen Kontexten also womöglich anders motiviert ist als in mündlichkeitsnahen Kontexten, dann kann eine mit Blick auf das gesprochene Deutsch formulierte Hypothese nur bedingt an einem schriftsprachlichen Korpus wie dem hier untersuchten Zeitschriftenkorpus überprüft werden, was ein besonderes Augenmerk auf eine umsichtige Wahl des Korpus lenkt. Um dieser Überlegung Rechnung zu tragen, bietet sich für Folgestudien an, die in einem solchen Korpus enthaltenen Textsorten (bspw. Kommentare, Leserbriefe und sonstige meinungsäußernde Texte) nach Möglichkeit getrennt voneinander zu analysieren, um so vielleicht zu einem differenzierteren Bild ohne Einfluss von potenziellen Störfaktoren zu gelangen. Durch dieses Vorgehen könnte sich in Anbetracht der zugrunde liegenden Hypothesen H1 und H2 ein direkteres Abbild ergeben als das in dieser Untersuchung ermittelte, bei der aufgrund fehlender textexterner Metadaten vonseiten der Korpusdokumentation keine textsortenentsprechende Trennung möglich \mbox{gewesen ist}.

\section{Fazit}

Die korpuslinguistische Studie startete mit der Beobachtung, dass es im Deutschen eine Vielzahl von Intensivierern gibt, woran sich die Frage anschließt, welche Faktoren ursächlich für diesen Variantenreichtum sind. Auf Grundlage der theoretischen Literatur lassen sich in diesem Kontext zwei Hypothesen nach \citet[2]{Hauschild.1899} und \citet[274]{Baumgarten.1908} ableiten, deren Gültigkeit mithilfe einer diachronen Frequenzanalyse auf den Grund gegangen wurde. Den Autoren nach geht die zunehmende Verwendung eines Intensivierers mit einem Verlust dessen expressiven Charakters einher, wobei neben der Frequenz als vielfach postuliertem Faktor erstmals auch die Persistenz als ausschlaggebend in Erwägung gezogen wurde (vgl. H1). Aus dem daraus resultierenden Expressivitätsverlust ergebe sich ein konstanter Bedarf an neuen Ausdrücken, die es Sprechern ermöglichen, zukünftig expressiv zu bleiben (vgl. H2). Vor diesem Hintergrund wurde angenommen, dass eine rückläufige Entwicklung von Intensivierern deren nachlassenden Expressivität geschuldet ist und das Zutreffen von H1 erhärtet. Neue Ausdrücke im Korpus wurden darüber hinaus als Evidenz für die Korrektheit von H2 herangezogen. Um die umrissenen Hypothesen zu testen, wurden für ausgewählte Intensivierer umfassende Korpusdaten erhoben.\footnote{An dieser Stelle sei erneut darauf hingewiesen, dass sich aus der Ableitung synchroner und diachroner Häufigkeitsverteilungen keine direkten Rückschlüsse auf den Expressivitätsgrad der Ausdrücke ziehen lassen. Für eine solche Bewertung sind experimentelle Untersuchungen unumgänglich. Dessen ungeachtet lassen sich -- wie in dieser Studie geschehen -- auf Grundlage der Ergebnisse der Korpusuntersuchung durchaus Annahmen über den Einfluss von Expressivität auf die Entwicklung eines Ausdrucks im Korpus formulieren, die in nachfolgenden Experimenten weiter überprüft werden können.}
\par Von den Ergebnissen ausgehend kann der Gebrauch des deskriptiven Intensivierers \textit{sehr} als eine Art Delle mit absinkendem Beginn und ansteigendem Ende beschrieben werden. Die noch andauernde Zunahme und das stabile Vorkommen erlauben die Interpretation, dass \textit{sehr} trotz temporärer Einbußen in den Neunzigerjahren auch weiterhin der stabilste Intensivierer des Deutschen ist und es keinen Grund zur Annahme einer Verdrängung durch andere Intensivierer gibt. Dies wurde mit seiner neutralen Gradbedeutung begründet, die dazu führt, dass der Gebrauch von \textit{sehr} ebenso in konservativen Textsorten und Registern vollkommen unproblematisch ist. Wie in Abschnitt \ref{Intensivierertypen} erörtert, hat \textit{sehr} seine ursprüngliche lexikalische Semantik unterdessen vollständig verloren, was ihm den Weg in die Rolle als rein skalarer Intensivierer geebnet hat. Infolgedessen ist das hohe Vorkommen von \textit{sehr} im Korpus insbesondere auf seine semantische Neutralität sowie seinen entsprechend hohen Lexikalisierungsgrad zurückzuführen. Dies stützt ferner auch die oben dargelegte Mutmaßung, nach der ein Verlust von Expressivität nicht notwendig eine Frequenzabnahme nach sich zieht, sodass dem Anschein nach wenig(er) expressive Intensivierer wie \textit{sehr} oder \textit{wahnsinnig} im Kontext einer sprachlich als konservativ geltenden Zeitschrift weniger stigmatisiert sind und entsprechend häufiger eingesetzt werden.

Im Kontext der expressiven Intensivierer deuten die diachronen Häufigkeitsverteilungen darauf hin, dass deren Verwendung im Korpus über den Untersuchungszeitraum hinweg linear angestiegen ist, wofür verschiedene Erklärungen in Betracht gezogen wurden. Bspw. besteht die Möglichkeit, dass sich der Schreibstil in den Wochenzeitschriften über die Zeit verändert hat, wodurch der Gebrauch expressiver Ausdrücke begünstigt worden ist. Außerdem könnte es zu einer stärkeren Gewichtung zum Vorteil konzeptionell mündlich orientierter Textsorten wie Interviews, Leserbriefen oder Kommentaren gekommen sein. Über diese Vermutung sind aufgrund unzureichender textexterner Metadaten und fehlender Textsortendokumentation vonseiten der Verlage jedoch keine harten Aussagen möglich. Demzufolge bleibt die Rolle der Registerhaftigkeit der Stellen, an denen die Intensivierer in den Texten auftraten, auch über die Untersuchung hinaus ungeklärt. Nichtsdestotrotz liefern Recherchen auf den frei zugänglichen Archiv-Seiten des \textsc{Spiegels}\footnote{\citet[]{Spiegel.2019}: \textit{https://www.spiegel.de/spiegel/print/index-2019.html} [\mbox{letzter Aufruf: 28.06.2021]}.} erste Indizien für diese Annahme, sind doch die neueren Ausgaben hinsichtlich der Sparten und der damit eingebundenen Textsorten merklich nuancierter als die älteren. Wie \mbox{\citet[73ff.\,94]{Just.1967}} in Bezug auf die inhaltliche Entwicklung des \textsc{Spiegels} berichtet, ist es in den ersten beiden Dekaden zu einem starken Ausbau von Meldungen, Kolumnen und Leserbriefen gekommen, die bekanntlich vornehmlich zur subjektiven Stellungnahme dienen: So hat sich der Anfangswert bei letzteren von ca. 0,6 Seiten ($\hateq$ 2,70\,\% des redaktionellen Teils) in 1947 auf ca. 3,7 Seiten \mbox{($\hateq$ 5,40\,\%} des redaktionellen Teils) in 1966 gesteigert \mbox{\citep[vgl.][76]{Just.1967}}. Daher ist denkbar, dass der Wandel zu mündlich orientierten Textsorten hin den Anstieg der Intensivierer in der Zeitschrift gefördert hat. Einen weiteren Hinweis auf zeitschrifteninterne Umstrukturierungen gibt \mbox{\citet[116]{Kuhn.1987}}, nach dem Printmedien über die Jahre einen verstärkten Fokus auf Anzeigenwerbung gerichtet haben: \glqq{}Der Aufschwung der Werbung in Deutschland ist unübersehbar, hat sich doch die Zahl der Anzeigeseiten in Publikumszeitschriften seit 1970 beinahe verdoppelt (von 85\,000 im Jahre 1970 auf 160\,000 im Jahr 1985)\grqq{}. Auch wenn eine Verdopplung der Werbeseiten beim \textsc{Spiegel} gemäß der von \citet[]{Kuhn.1987} angegebenen Zahlen (5\,485 Seiten im Jahr 1970 und 6\,240 Seiten in 1986) etwas zu hoch gegriffen ist, ist dennoch ein klarer Anstieg zu erkennen. Demnach scheint das Nachrichtenmagazin vermehrt auf Anzeigenschaltung als Werbestrategie zur Markenrepräsentation zu setzen, wie bspw. \citet[31f.,\,115]{Just.1967} hervorhebt. Ähnliches trifft womöglich auf die \textsc{Zeit} zu, bei der das höhere Anzeigenvolumen durch eine stärkere Auflage und eine Zunahme im Umfang bedingt ist \mbox{\citep[vgl.][260,\,422]{Jansen.2006}}. Bedenkt man, dass der Gebrauch expressiver Sprache ein gerne eingesetztes Mittel des Marketings ist, um das Interesse potenzieller Käufer zu wecken und Kaufanreize zu schaffen \citep[vgl.][97]{Suscinskij.1985}, ist durchaus denkbar, dass der Anstieg im Bereich der Anzeigenwerbung mitunter Anteil an der stetigen Zunahme der Intensivierer im Korpus hatte. Eine weitere Erklärung für den Anstieg tangierte die Definition der Auswirkungen von Expressivitätsverlust. Während sich dieser gemäß den Expressivitätsmaximen nach \mbox{\citet[]{Kirschbaum.2002b}} (vgl. Abschnitt \ref{Auslöser}) in der gesprochenen Sprache in Form eines Frequenzrückgangs ausdrückt, ist vorstellbar, dass er sich in der geschriebenen Sprache hingegen vielmehr in einem Anstieg der Frequenz niederschlägt. Denn je unkritischer der Gebrauch eines Intensivierers ist, desto besser ist er für schriftsprachliche Kontexte, allen voran für die im Mittelpunkt dieser Untersuchung stehenden Zeitungssprache, geeignet. In diesem Sinne schlägt sich der Verlust von Expressivität in Abhängigkeit des Realisierungsmediums nach \mbox{\citet[]{Koch.1996,Koch.2007}} mutmaßlich in unterschiedlicher, gar divergierender Weise nieder. Von den Ergebnissen der Untersuchung ausgehend besteht also die Möglichkeit, dass die hochfrequenten Intensivierer über die Jahre hinweg an Expressivität \mbox{eingebüßt haben}.

Um dem fraglichen Expressivitätsverlust nachzugehen, wurde der Untersuchungszeitraum in drei Zeitsegmente gegliedert, in denen die expressiven Intensivierer genauer betrachtet wurden. Die Begutachtung erfolgte mit Blick auf die Hypothese H1, weshalb an dieser Stelle die beiden Faktoren Erstvorkommen (vgl. H1-Temp) und Frequenz (vgl. H1-Freq) im Fokus standen. In diesem Kontext zeigte sich, dass die Gruppenmitglieder neben ähnlichen Frequenzen ebenso semantische Gemeinsamkeiten aufweisen, die eine mögliche Erklärung für die Häufigkeiten im Korpus bieten. So hat die Semantik dem Anschein nach insofern zur Gruppenbildung beigetragen, als sie den Gebrauch der Ausdrücke in einigen Fällen begünstigt (vgl. Cluster der hochfrequenten Intensivierer), ihn in anderen Fällen allerdings behindert (vgl. Cluster der niedrigfrequenten Ausdrücke) zu haben scheint. Vor diesem Hintergrund wurde angenommen, dass die Verwendung insbesondere der Intensivierer mit klar erkennbarer und negativer Semantik (z.\,B. \textit{sau} und \textit{scheiß}) in der Zeitungssprache wesentlich kontroverser ist als der Gebrauch der Intensivierer des zweiten Clusters. Dies deutet auf einen semantischen Filter hin, der den Gebrauch ebensolcher offensiven Ausdrücke in der Schriftsprache im Allgemeinen, vor allem aber in konservativen Textsorten beschränkt. In der hochfrequenten Gruppe weisen die Mitglieder der \glqq{}schrecklichen\grqq{} Intensivierergruppe zwar gleichfalls eine transparente Semantik auf, jedoch ist diese in intensivierender Funktion aufgrund von Abnutzung eventuell nicht weiter von Bedeutung \citep[vgl.][15f.]{Nouwen.2020}. Auch hier könnten weiterführende experimentelle Untersuchungen Aufschluss darüber geben, ob und inwieweit die zugrunde liegende Semantik der Ausdrücke auf die Interpretation des Intensivausdrucks einwirkt. Wie sich die Frequenz der expressiven Intensivierer im Korpus in den nächsten Jahren entwickeln wird, kann zum jetzigen Zeitpunkt naturgemäß nicht beantwortet werden.

Zusammengenommen steigt gegenwärtig der Gebrauch beinahe aller untersuchten Intensivierer an. Dabei können die rezent in Erscheinung getretenen expressiven Intensivierer \textit{hammer}, \textit{krass} und \textit{mega} als Evidenz für die auf \citet[274]{Baumgarten.1908} zurückgehende Hypothese H2 herangezogen werden, nach der es einen permanenten Bedarf an neuen, unverbrauchten Ausdrücken gibt. Wie oben beschrieben, sind diese Ausdrücke erstmals im Zeitraum um 1970 im Untersuchungskorpus aufgetreten. Der Ausdruck \textit{mega} wird im Rahmen der von \citet[345]{Ruf.1996} durchgeführten Korpusstudie zu augmentativen Lehnpräfixen ebenso nicht nur als neu, sondern auch als überaus produktiv beschrieben: \glqq{}Es ist \textelp{} anzunehmen, daß \textit{Mega-} den Höhepunkt seiner Produktivität 1996 noch nicht erreicht hat.\grqq{} Angesichts der im Korpus gegenwärtig stark ausgeprägten Zunahme des intensivierenden \textit{mega} teile ich diese Vermutung. So ist zu erwarten, dass die positive Entwicklung weiterhin andauert und \textit{mega} auch ein Vierteljahrhundert später sein Frequenzmaximum noch nicht erreicht hat. Daneben konnte im Korpus ein Intensivierer identifiziert werden, dessen Gebrauch im Laufe der Zeit stärker abgenommen hat und infolgedessen zweifellos einen Grund zur Annahme des von \citet[2]{Hauschild.1899} in H1 postulierten Expressivitätsverlust bietet, nämlich \textit{derb}. Hier geben die relativen Häufigkeiten über den Untersuchungszeitraum von nahezu 70 Jahren zu verstehen, dass der Ausdruck im Vergleich zu den restlichen Intensivierern erheblich seltener gebraucht wird, sodass er im Korpus inzwischen fast gar nicht mehr in intensivierender Funktion vorkommt. Aufgrund der Vielfalt und der fortlaufenden Entwicklung der im Deutschen verwendeten Intensivierer erscheint es realistisch, dass das intensivierende \textit{derb} zukünftig in der mündlichen wie auch der schriftlichen Kommunikation obsolet sein wird, auch wenn sein tatsächlicher Expressivitätsgrad bis hierhin noch ungeklärt ist. Darüber hinaus besteht die Möglichkeit, dass auch hochfrequente Intensivierer wie \textit{verdammt} oder \textit{wahnsinnig} an Expressivität verloren haben, gleichwohl deren Gebrauch über die Jahre hinweg stetig angestiegen ist und er im Sinne der Expressivitätsmaximen nach \citet[]{Kirschbaum.2002b} zunächst keine Hinweise auf einen Expressivitätsverlust gibt. Wie oben diskutiert, schlägt er sich im schriftsprachlichen Kontext gegebenenfalls quantitativ jedoch mehr in einem Anstieg als einem Rückgang nieder, dessen Ursache in einer besseren Eignung, einer größeren Flexibilität sowie einer entsprechend höheren Frequenz begründet liegt. Auch weist diese Definition von Expressivitätsverlust auf medienspezifische Unterschiede hin, sodass sich der Rückgang von Expressivität in Abhängigkeit des Kommunikationsmediums nach \citet[]{Koch.1996,Koch.2007} auf unterschiedliche, vielleicht sogar divergierende Weise ausdrückt.

Die eingangs beschriebenen Bedenken, nach denen es schwierig sein kann, im zugrunde liegenden Korpus genügend Belege für typischerweise mündlichkeitsnahe Phänomene zu finden, haben sich nicht bestätigt. Somit lässt sich festhalten, dass das Vorkommen von expressiven Intensivierern auch in professionell redigierten Textformen wie Zeitschriften trotz in manchen Fällen minderer Ergiebigkeit erhoben werden kann. Durch die dadurch ermittelten Entwicklungstendenzen konnten die Ausdrücke wie intendiert quantitativ eingeordnet werden. Die beobachtete Frequenzzunahme der expressiven Intensivierer kann ein Indikator sowohl für einen Expressivitätsverlust als auch einen fortgeschrittenen Sprachwandel sein, lasse sich \glqq{}anhand der Ausbreitung indirekt schließen, dass sich das, was häufig ist, bewährt hat.\grqq{} (\citealt[70]{Perkuhn.2012}). Auf die Ergebnisse der Untersuchung übertragen, ist daher folgender Schluss möglich: Treten über die Zeit vermehrt Belege für die expressiven Intensivierer im Korpus auf, so ist nicht ausgeschlossen, dass die Ausdrücke in der mündlichen Kommunikation voll etabliert sind und der Wandel sukzessive auf den schriftlichen Sprachgebrauch übergegriffen hat.\footnote{Nichtsdestotrotz sollte berücksichtigt werden, dass der schriftliche Sprachgebrauch deutlich konservativer ist als der mündliche. Insofern können Belege im Korpus den Eindruck erwecken, ein Ausdruck sei noch gebräuchlich, obwohl er in der gesprochenen Sprache längst obsolet ist.} Während die Hypothese H2 nach \citet[274]{Baumgarten.1908} den Ergebnissen der Korpusanalyse nach de facto gültig zu sein scheint, ist eine Beantwortung der Hypothese H1 nach \citet[2]{Hauschild.1899} durch eine korpusgestützte Häufigkeitsanalyse lediglich marginal möglich. Dennoch liefern die Ergebnisse erste empirische Evidenz für einen generellen Anstieg von Intensivierern in den Wochenzeitschriften \textsc{Der Spiegel} und \textsc{Die Zeit} und können als breite Datengrundlage für nachfolgende experimentelle Untersuchungen dienen.

Eine solche weiterführende Einordnung wird in den Kapiteln \ref{Experiment0} und \ref{Experimente} dieses Buchs vorgenommen. Wie zu Beginn zur Sprache gebracht, wurden ebensolche experimentellen Verfahren meines Wissens bis dato nicht angewandt, weswegen die vorliegende Untersuchung einen ersten Schritt in diese Richtung macht. So knüpft das in Kapitel \ref{Experiment0} vorgestellte Experiment an die Erkenntnisse und Schlussfolgerungen der Korpusstudie an und überprüft das Zutreffen der nach wie vor zur Diskussion stehenden Hypothesen H1-Freq und H1-Temp nach \citet[2]{Hauschild.1899} und \citet[59]{Tobler.1868}. Im Zuge der Untersuchung wird der Frage nach dem Expressivitätsgrad der Intensivierer und demnach der Gültigkeit von H1 unter Berücksichtigung der Faktoren Frequenz (vgl. H1-Freq) und Erstvorkommen (vgl. H1-Temp) weiter auf den Grund gegangen. Von elementarer Bedeutung ist die Beantwortung der Fragen, ob die intraklassischen Unterschiede der Korpusstudie in gleicher Weise bei experimentellen Untersuchungen zutage treten und ob sich die Ergebnisse beider empirischen linguistischen Herangehensweisen decken oder divergieren. Um den Expressivitätsgrad der Ausdrücke zu erheben, werden Intensivausdrücke mit den in der Korpusstudie berücksichtigten Intensivierern im Experiment von einer umfangreichen Personenstichprobe auf einer Skala relativ zueinander in Beziehung gesetzt.
